\documentclass[../libro.tex]{subfiles}

\begin{document}

\ifSubfilesClassLoaded{\appendix\chapter{Afterstories}\clearpage}{}

\section{线性方程组与秩}

本节, 我们讨论线性方程组与阵的秩.

回想, 线性方程组是形如 \(A X = B\) 的阵等式,
其中, \(A\), \(B\) 是 \(m \times n\), \(m \times 1\)~阵,
且 \(X\) 是未知的 \(n \times 1\)~阵.
若 \(n \times 1\)~阵 \(C\) 适合 \(A C = B\),
则 \(C\) 是 \(A X = B\) 的一个解.

在第一章, 节~\malneprasekcio{27}, 我们有

\begin{theorem}
    设 \(A\) 是 \(m \times n\)~阵.
    设 \(B\) 是 \(m \times 1\)~阵.
    作一个 \(m \times (n+1)\)~阵 \(G\),
    其中,
    \begin{align*}
        [G]_{i,j}
        = \begin{cases}
              [A]_{i,j}, & j \leq n;  \\
              [B]_{i,1}, & j = n + 1.
          \end{cases}
    \end{align*}
    (通俗地, 在 \(A\)~的最后一列的右侧加入一列 \(B\),
    得到尺寸较大的阵 \(G\).)
    设
    \(A\) 有一个行列式非零的 \(r\)~级子阵,
    但没有行列式非零的 \(r+1\)~级子阵.

    (1)
    若 \(G\)~也没有行列式非零的 \(r+1\)~级子阵,
    则 \(AX = B\) 有解.
    进一步地, 若 \(r < n\),
    则 \(AX = B\) 的解不唯一;
    若 \(r = n\),
    则 \(AX = B\) 的解唯一.

    (2)
    若 \(G\)~有一个行列式非零的 \(r+1\)~级子阵,
    则 \(AX = B\) 无解.
\end{theorem}

本节, 我们推广此事, 且定义与讨论阵的秩.

我们说, 行列式是方阵的一个重要的属性.
不过, 不是方阵的阵没有行列式,
故行列式不是不是方阵的阵的属性.
但是, 秩是每个阵都有的一个重要的属性.

\begin{theorem}
    设 \(A\) 是 \(m \times n\)~阵.
    设 \(B\) 是 \(m \times t\)~阵.
    作 \(m \times (n + t)\)~阵 \(G\), 其中,
    \begin{align*}
        [G]_{i,j} =
        \begin{cases}
            [A]_{i,j},   & j \leq n; \\
            [B]_{i,j-n}, & j > n.
        \end{cases}
    \end{align*}
    设
    \(A\) 有一个行列式非零的 \(r\)~级子阵,
    但没有行列式非零的 \(r+1\)~级子阵.

    (1)
    设存在 \(n \times t\)~阵 \(C\), 使 \(A C = B\).
    则 \(G\) 没有行列式非零的 \(r+1\)~级子阵.

    (2)
    反过来, 设 \(G\) 没有行列式非零的 \(r+1\)~级子阵.
    则存在 \(n \times t\)~阵 \(C\), 使 \(A C = B\).
    进一步地, 若 \(r < n\), 则这样的 \(C\) 不唯一;
    若 \(r = n\), 则这样的 \(C\) 唯一.
\end{theorem}

注意, 若我们取 \(t = 1\), 则这就是上个定理.

\begin{proof}
    (1)
    设 \(n \times t\)~阵 \(C\) 使 \(A C = B\).

    若 \(G\) 没有 \(r + 1\)~级子阵,
    则 \(G\) 没有行列式非零的 \(r + 1\)~级子阵.

    下设 \(G\) 有 \(r + 1\)~级子阵.
    则 \(r + 1 \leq m\),
    且 \(r + 1 \leq n + t\).

    作 \(n \times (n + t)\)~阵 \(H\),
    其中,
    \begin{align*}
        [H]_{i,j} =
        \begin{cases}
            [I_n]_{i,j}, & j \leq n; \\
            [C]_{i,j-n}, & j > n.
        \end{cases}
    \end{align*}
    则 \(AH\) 是 \(n \times (n + t)\) 阵.
    当 \(j \leq n\) 时,
    \begin{align*}
        [A H]_{i,j}
        = {} &
        \sum_{1 \leq \ell \leq n}
        {[A]_{i,\ell} [H]_{\ell,j}}
        \\
        = {} &
        \sum_{1 \leq \ell \leq n}
        {[A]_{i,\ell} [I_n]_{\ell,j}}
        \\
        = {} &
        [A I_n]_{i,j}
        \\
        = {} &
        [A]_{i,j}
        \\
        = {} &
        [G]_{i,j};
    \end{align*}
    当 \(j > n\) 时,
    \begin{align*}
        [A H]_{i,j}
        = {} &
        \sum_{1 \leq \ell \leq n}
        {[A]_{i,\ell} [H]_{\ell,j}}
        \\
        = {} &
        \sum_{1 \leq \ell \leq n}
        {[A]_{i,\ell} [C]_{\ell,j-n}}
        \\
        = {} &
        [A C]_{i,j-n}
        \\
        = {} &
        [B]_{i,j-n}
        \\
        = {} &
        [G]_{i,j}.
    \end{align*}
    则 \(AH = G\).

    设 \(1 \leq i_1 < i_2 < \dots < i_{r+1} \leq m\).
    设 \(1 \leq j_1 < j_2 < \dots < j_{r+1} \leq n+t\).
    作 \(G\) 的一个 \(r + 1\)~级子阵
    \begin{align*}
        G' = G\binom{i_1,i_2,\dots,i_{r+1}}{j_1,j_2,\dots,j_{r+1}}.
    \end{align*}
    我们证明, \(\det {(G')} = 0\).

    考虑
    \begin{align*}
        A' = A\binom{i_1,i_2,\dots,i_{r+1}}{1,2,\dots,n},
        \quad
        H' = H\binom{1,2,\dots,n}{j_1,j_2,\dots,j_{r+1}}.
    \end{align*}
    则 \([A']_{u,v} = [A]_{i_u,v}\),
    且 \([H']_{u,v} = [H]_{u,j_v}\).
    则 \(A' H'\) 是 \(r + 1\)~级阵.
    注意, \([G']_{u,v} = [G]_{i_u,j_v}\).
    注意,
    \begin{align*}
        [A' H']_{u,v}
        = {} &
        \sum_{1 \leq \ell \leq n}
        {[A']_{u,\ell} [H']_{\ell,v}}
        \\
        = {} &
        \sum_{1 \leq \ell \leq n}
        {[A]_{i_u,\ell} [H]_{\ell,j_v}}
        \\
        = {} &
        [A H]_{i_u,j_v}
        \\
        = {} &
        [G]_{i_u,j_v}
        \\
        = {} &
        [G']_{u,v}.
    \end{align*}
    则 \(A' H' = G'\).

    设 \(A'\) 的列
    \(1\), \(2\), \(\dots\), \(n\)
    分别是 \(a_1\), \(a_2\), \(\dots\), \(a_n\).
    设 \(H'\) 的列
    \(1\), \(2\), \(\dots\), \(r+1\)
    分别是 \(h_1\), \(h_2\), \(\dots\), \(h_{r+1}\).
    则
    \begin{align*}
        G'
        = {} &
        A' H'
        \\
        = {} &
        [A' h_1, A' h_2, \dots, A' h_{r+1}]
        \\
        = {} &
        \left[
        \sum_{k_1 = 1}^{n} {[h_1]_{k_1,1} a_{k_1}},
        \sum_{k_2 = 1}^{n} {[h_2]_{k_2,1} a_{k_2}},
        \dots,
        \sum_{k_{r+1} = 1}^{n} {[h_{r+1}]_{k_{r+1},1} a_{k_{r+1}}}
        \right]
        \\
        = {} &
        \left[
        \sum_{k_1 = 1}^{n} {[H']_{k_1,1} a_{k_1}},
        \sum_{k_2 = 1}^{n} {[H']_{k_2,2} a_{k_2}},
        \dots,
        \sum_{k_{r+1} = 1}^{n} {[H']_{k_{r+1},r+1} a_{k_{r+1}}}
        \right].
    \end{align*}
    由行列式的多线性,
    \begin{align*}
             &
        \det {(G')}
        \\
        = {} &
        \det {
            \left[
                \sum_{k_1 = 1}^{n} [H']_{k_1,1} a_{k_1},
                \sum_{k_2 = 1}^{n} [H']_{k_2,2} a_{k_2},
                \dots,
                \sum_{k_{r+1} = 1}^{n} [H']_{k_{r+1},r+1} a_{k_{r+1}}
                \right]
        }
        \\
        = {} &
        \sum_{k_1 = 1}^{n}
        [H']_{k_1,1}
        \det {
            \left[
                a_{k_1},
                \sum_{k_2 = 1}^{n} [H']_{k_2,2} a_{k_2},
                \dots,
                \sum_{k_{r+1} = 1}^{n} [H']_{k_{r+1},r+1} a_{k_{r+1}}
                \right]
        }
        \\
        = {} &
        \sum_{k_1 = 1}^{n}
        \sum_{k_2 = 1}^{n}
        [H']_{k_1,1}
        [H']_{k_2,2}
        \det {
            \left[
                a_{k_1},
                a_{k_2},
                \dots,
                \sum_{k_{r+1} = 1}^{n} [H']_{k_{r+1},r+1} a_{k_{r+1}}
                \right]
        }
        \\
        = {} &
        \dots
        \\
        = {} &
        \sum_{k_1 = 1}^{n}
        \sum_{k_2 = 1}^{n}
        \dots
        \sum_{k_{r+1} = 1}^{n}
        [H']_{k_1,1}
        [H']_{k_2,2}
        \dots
        [H']_{k_{r+1},r+1}
        \det {[a_{k_1}, a_{k_2}, \dots, a_{k_{r+1}}]}
        \\
        = {} &
        \sum_{1 \leq k_1, k_2, \dots, k_{r+1} \leq n}
        [H']_{k_1,1}
        [H']_{k_2,2}
        \dots
        [H']_{k_{r+1},r+1}
        \det {[a_{k_1}, a_{k_2}, \dots, a_{k_{r+1}}]}.
    \end{align*}

    若在 \(k_1\), \(k_2\), \(\dots\), \(k_{r+1}\),
    有二个数相同,
    则 \([a_{k_1}, a_{k_2}, \dots, a_{k_{r+1}}]\)
    有相同的二列.
    由行列式的交错性,
    \(\det {[a_{k_1}, a_{k_2}, \dots, a_{k_{r+1}}]} = 0\).

    若 \(k_1\), \(k_2\), \(\dots\), \(k_{r+1}\)
    互不相同
    (注意, 此时, \(n \geq r+1\)),
    我们适当地交换
    \([a_{k_1}, a_{k_2}, \dots, a_{k_{r+1}}]\)
    的列的次序,
    以变它为 \(A'\) 的一个 \(r+1\)~级子阵,
    其当然是 \(A\) 的一个 \(r+1\)~级子阵.
    由行列式的反称性,
    \begin{align*}
        \det {[a_{k_1}, a_{k_2}, \dots, a_{k_{r+1}}]}
        = {} &
        {\pm
                \det {\left(
                    A'\binom{1,2,\dots,r+1}{k_1,k_2,\dots,k_{r+1}}
                    \right)}}
        \\
        = {} &
        {\pm
                \det {\left(
                    A\binom{i_1,i_2,\dots,i_{r+1}}{k_1,k_2,\dots,k_{r+1}}
                    \right)}}
        \\
        = {} &
        0.
    \end{align*}
    于是, \(\det {(G')}\) 是一些 \(0\) 的和.
    则 \(\det {(G')} = 0\).

    (2)
    设 \(G\) 没有行列式非零的 \(r + 1\)~级子阵.

    设 \(B\) 的列
    \(1\), \(2\), \(\dots\), \(t\)
    分别是 \(b_1\), \(b_2\), \(\dots\), \(b_t\).
    作 \(t\)~个 \(m \times (n + 1)\) 阵 \(G_v\)
    (\(v = 1\), \(2\), \(\dots\), \(t\)):
    \begin{align*}
        [G_v]_{i,j} =
        \begin{cases}
            [A]_{i,j},               & j \leq n;  \\
            [b_v]_{i,1} = [B]_{i,v}, & j = n + 1.
        \end{cases}
    \end{align*}
    不难看出, \(G_v\) 是 \(G\) 的子阵,
    故 \(G_v\) 也没有行列式非零的 \(r + 1\)~级子阵.
    则由上个定理, 存在 \(n \times 1\)~阵 \(c_v\),
    使 \(A c_v = b_v\).

    作 \(n \times t\)~阵
    \(C = [c_1, c_2, \dots, c_t]\).
    则 \(A C\) 是 \(m \times t\)~阵,
    且
    \begin{align*}
        [A C]_{i,j}
        = {} &
        \sum_{1 \leq \ell \leq n}
        {[A]_{i,\ell} [C]_{\ell,j}}
        \\
        = {} &
        \sum_{1 \leq \ell \leq n}
        {[A]_{i,\ell} [c_j]_{\ell,1}}
        \\
        = {} &
        [A c_j]_{i,1}
        \\
        = {} &
        [b_j]_{i,1}
        \\
        = {} &
        [B]_{i,j}.
    \end{align*}
    则 \(A C = B\).

    注意, 若 \(r < n\),
    则, 由上个定理,
    存在 \(n \times 1\)~阵 \(c'_1\),
    使 \(c'_1 \neq c_1\),
    且 \(A c'_1 = b_1\).
    再作 \(C' = [c'_1, c_2, \dots, c_t]\).
    则我们可类似地证明, \(A C' = B\).
    不过, \(C'\) 当然不等于 \(C\),
    因为 \(C'\) 与 \(C\) 的列~\(1\) 不相等.
    于是, 若 \(r < n\), 则适合 \(A C = B\)
    的 \(C\) 不唯一.

    最后, 注意, 若 \(r = n\),
    则 \(m \times n\)~阵 \(A\)
    有行列式非零的 \(n\)~级子阵.
    设 \(n \times t\)~阵 \(P\), \(Q\) 适合,
    \(A P = B\), 且 \(A Q = B\).
    则 \(A P = A Q\).
    由消去律, \(P = Q\)
    (见 ``消去律'').
    于是, 若 \(r = n\), 则适合 \(A C = B\)
    的 \(C\) 唯一.
\end{proof}

由此可见, 对一个阵,
``有一个行列式非零的 \(r\)~级子阵,
但没有行列式非零的 \(r+1\)~级子阵''
的 \(r\) 或许是重要的属性.
我们定义

\begin{definition}
    设 \(A\) 是 \(m \times n\)~阵.
    则存在唯一的非负整数 \(r\),
    使
    (见第一章, 节~\malneprasekcio{25}):

    (1)
    \(A\) 有一个行列式非零的 \(r\)~级子阵;

    (2)
    \(A\) 没有行列式非零的 \(r+1\)~级子阵.

    我们定义 \(A\) 的\emph{秩},
    \(\operatorname{rank} {(A)}\),
    为 \(r\).

    注意, 若 \(A = 0\), 则 \(r = 0\)
    (因为我们约定, ``\(0\)~级阵'' 的行列式为 \(1\)).
\end{definition}

回想, 若 \(A\) 是 \(m \times n\)~阵,
且 \(B\) 是 \(m \times t\)~阵,
则 \([A, B]\) 是 \(m \times (n + t)\)~阵 \(G\),
其中,
\begin{align*}
    [G]_{i,j} =
    \begin{cases}
        [A]_{i,j},   & j \leq n; \\
        [B]_{i,j-n}, & j > n.
    \end{cases}
\end{align*}
则我们可如此改写上个定理的 (1):

\begin{theorem}
    设 \(A\) 是 \(m \times n\)~阵.
    设 \(C\) 是 \(n \times t\)~阵.
    则
    \begin{align*}
        \operatorname{rank} {([A, AC])}
        = \operatorname{rank} {(A)}.
    \end{align*}
\end{theorem}

\begin{proof}
    设 \(r = \operatorname{rank} {(A)}\).
    则 \(A\) 有行列式非零的 \(r\)~级子阵,
    且 \(A\) 没有行列式非零的 \(r + 1\)~级子阵.
    \(A\) 是 \([A, AC]\) 的子阵,
    故 \([A, AC]\) 也有行列式非零的 \(r\)~级子阵.
    由上个定理的 (1),
    \([A, AC]\) 没有行列式非零的 \(r + 1\)~级子阵.
    则 \(\operatorname{rank} {([A, AC])} = r\).
\end{proof}

% \begin{proof}
%     设 \(r = \operatorname{rank} {(A)}\).
%     由上个定理的 (1),
%     \(\operatorname{rank} {([A, AC])} \leq r\).
%     另一方面, 既然 \(A\) 是 \([A, AC]\) 的子阵,
%     则 \(\operatorname{rank} {([A, AC])} \geq r\).
% \end{proof}

\begin{theorem}
    设 \(A\) 是 \(m \times n\)~阵.
    则 \(\operatorname{rank} {(A)}
    = \operatorname{rank} {(A^{\mathrm{T}})}\).
\end{theorem}

\begin{proof}
    注意, \(A^{\mathrm{T}}\) 的子阵是 \(A\) 的某子阵的转置,
    且 \((A^{\mathrm{T}})^{\mathrm{T}} = A\)
    的子阵也是 \(A^{\mathrm{T}}\) 的某子阵的转置.
    于是, 若 \(A\) 有行列式非零的 \(r\)~级子阵,
    则 \(A^{\mathrm{T}}\) 也有行列式非零的 \(r\)~级子阵,
    且若 \(A\) 没有行列式非零的 \(r + 1\)~级子阵,
    则 \(A^{\mathrm{T}}\) 也没有行列式非零的 \(r + 1\)~级子阵.
\end{proof}

\begin{theorem}
    设 \(A\) 是 \(m \times n\)~阵.
    设 \(C\) 是 \(n \times t\)~阵.
    则
    \begin{align*}
        \operatorname{rank} {(AC)}
        \leq \operatorname{rank} {(A)}.
    \end{align*}
\end{theorem}

\begin{proof}[证~1]
    设 \(r = \operatorname{rank} {(A)}\).
    记 \(B = A C\).
    我们证明, \(B\) 没有行列式非零的 \(r + 1\)~级子阵.
    若 \(B\) 没有 \(r + 1\)~级子阵,
    则 \(B\) 没有行列式非零的 \(r + 1\)~级子阵.
    若 \(B\) 有 \(r + 1\)~级子阵,
    我们任取一个 \(r + 1\)~级子阵
    \begin{align*}
        B' =
        B\binom{i_1,i_2,\dots,i_{r+1}}{j_1,j_2,\dots,j_{r+1}},
    \end{align*}
    其中, \(1 \leq i_1 < i_2 < \dots i_{r+1} \leq m\),
    且 \(1 \leq j_1 < j_2 < \dots < j_{r+1} \leq t\),
    并证明, \(\det {(B')} = 0\).
    就像前面的讨论那样, 我们考虑
    \begin{align*}
        A' = A\binom{i_1,i_2,\dots,i_{r+1}}{1,2,\dots,n},
        \quad
        C' = C\binom{1,2,\dots,n}{j_1,j_2,\dots,j_{r+1}}.
    \end{align*}
    就像前面的讨论那样, 我们可验证,
    \(A' C' = B'\).
    于是, 就像前面的讨论那样, 我们可证明,
    \(\det {(B')} = 0\).
    请允许我留它们为您的习题.
\end{proof}

\begin{proof}[证~2]
    我们用已有的结论.
    注意, \(AC\) 是 \([A, AC]\) 的子阵,
    故 \(\operatorname{rank} {(AC)}
    \leq \operatorname{rank} {([A, AC])}\).
    则由前面的讨论,
    \begin{equation*}
        \operatorname{rank} {(AC)}
        \leq \operatorname{rank} {([A, AC])}
        = \operatorname{rank} {(A)}.
        \qedhere
    \end{equation*}
\end{proof}

\begin{theorem}
    设 \(A\) 是 \(m \times n\)~阵.
    设 \(C\) 是 \(n \times t\)~阵.
    则
    \begin{align*}
        \operatorname{rank} {(AC)}
         & \leq \operatorname{rank} {(A)}, \\
        \operatorname{rank} {(AC)}
         & \leq \operatorname{rank} {(C)}.
    \end{align*}
\end{theorem}

\begin{proof}
    不等式~1 已被证明.
    我们现在证明不等式~2.
    注意,
    \begin{align*}
        \operatorname{rank} {(AC)}
        = {}    &
        \operatorname{rank} {((AC)^{\mathrm{T}})}
        \\
        = {}    &
        \operatorname{rank} {(C^{\mathrm{T}} A^{\mathrm{T}})}
        \\
        \leq {} &
        \operatorname{rank} {(C^{\mathrm{T}})}
        \\
        = {}    &
        \operatorname{rank} {(C)}.
        \qedhere
    \end{align*}
\end{proof}

\begin{theorem}
    设 \(A\) 是 \(m \times n\)~阵.
    设 \(P\) 是 \(n\)~级阵, 且 \(\det {(P)} \neq 0\).
    设 \(Q\) 是 \(m\)~级阵, 且 \(\det {(Q)} \neq 0\).
    则
    \begin{align*}
        \operatorname{rank} {(QAP)}
        = \operatorname{rank} {(A)}.
    \end{align*}
\end{theorem}

\begin{proof}
    设 \(P_1 = (\det {(P)})^{-1}
    \operatorname{adj} {(P)}\),
    且 \(Q_1 = (\det {(Q)})^{-1}
    \operatorname{adj} {(Q)}\).
    则 \(P P_1 = I_n\),
    且 \(Q_1 Q = I_m\).
    则
    \begin{align*}
        \operatorname{rank} {(A)}
        \geq {} &
        \operatorname{rank} {(AP)}
        \\
        \geq {} &
        \operatorname{rank} {(QAP)}
        \\
        \geq {} &
        \operatorname{rank} {(QAP P_1)}
        \\
        \geq {} &
        \operatorname{rank} {(Q_1 QAP P_1)}
        \\
        = {}    &
        \operatorname{rank} {((Q_1 Q) A (P P_1))}
        \\
        = {}    &
        \operatorname{rank} {(I_m A I_n)}
        \\
        = {}    &
        \operatorname{rank} {(A)}.
    \end{align*}
    则
    \(\operatorname{rank} {(A)}
    \geq \operatorname{rank} {(Q A P)}
    \geq \operatorname{rank} {(A)}\).
\end{proof}

由此可见, 倍加不改变秩:
对 \(m \times n\)~阵 \(A\),
\begin{align*}
    \operatorname{rank} {(A E(n; p, q; s))}
     & = \operatorname{rank} {(A)}, \\
    \operatorname{rank} {(E(m; p', q'; s') A)}
     & = \operatorname{rank} {(A)}.
\end{align*}
(回想, \(E(n; p, q; s) E(n; p, q; -s) = I_n\),
且 \(E(m; p', q'; -s') E(m; p', q'; s') = I_m\),
若 \(p \neq q\), 且 \(p' \neq q'\).)

\begin{example}
    设 \(m \times n\)~阵 \(D\) 适合,
    若 \(i \neq j\), 则 \([D]_{i,j} = 0\).
    设存在 \(r\)~个正整数
    \(i_1\), \(i_2\), \(\dots\), \(i_r\)
    (\(1 \leq i_1 < i_2 < \dots < i_r\),
    且 \(i_r \leq m\), 且 \(i_r \leq n\)),
    使 \([D]_{i_k,i_k} \neq 0\)
    (\(k = 1\), \(\dots\), \(r\)),
    且 \(u\) 不等于 \(i_1\), \(\dots\), \(i_r\)
    的任何一个时, \([D]_{u,u} = 0\).
    则 \(\operatorname{rank} {(D)} = r\).

    首先, \(D\) 有一个行列式非零的 \(r\)~级子阵,
    因为
    \begin{align*}
        \det {\left(
            D\binom{i_1,i_2,\dots,i_r}{i_1,i_2,\dots,i_r}
            \right)}
        = {}    &
        \det {
            \begin{bmatrix}
                [D]_{i_1,i_1} & 0             & \cdots & 0             \\
                0             & [D]_{i_2,i_2} & \cdots & 0             \\
                \vdots        & \vdots        & \ddots & \vdots        \\
                0             & 0             & \cdots & [D]_{i_r,i_r} \\
            \end{bmatrix}
        }
        \\
        = {}    &
        [D]_{i_1,i_1} \dots [D]_{i_r,i_r}
        \\
        \neq {} &
        0.
    \end{align*}
    然后, \(D\) 没有行列式非零的 \(r + 1\)~级子阵.
    若 \(D\) 没有 \(r + 1\)~级子阵,
    则 \(D\) 当然没有行列式非零的 \(r + 1\)~级子阵.
    设 \(D\) 有 \(r + 1\)~级子阵.
    那么, 在 \(D\) 的任何 (不同的) \(r + 1\) 列,
    必存在一列, 其是零.
    (因为当 \(v\) 不等于 \(i_1\), \(\dots\), \(i_r\)
    的任何一个时,
    \([D]_{v,v} \neq 0\),
    且若 \(u \neq v\),
    则 \([D]_{u,v} \neq 0\).
    则当 \(v\) 不等于 \(i_1\), \(\dots\), \(i_r\)
    的任何一个时,
    \(D\) 的列~\(v\) 的元全为 \(0\).)
    则 \(D\) 的任何 \(r + 1\)~级子阵%
    必有一列, 其是零.
    则它的行列式为 \(0\).
    则 \(\operatorname{rank} {(D)} = r\).
\end{example}

最后, 我们讨论 \(n\)~级阵 (\(n \geq 2\)) \(A\)
的古伴 \(\operatorname{adj} {(A)}\) 的秩.

若 \(\operatorname{rank} {(A)} = n\),
则 \(\det {(A)} \neq 0\).
则 \(\det {(\operatorname{adj} {(A)})}
= (\det {(A)})^{n-1} \neq 0\).
则 \(\operatorname{rank} {(\operatorname{adj} {(A)})} = n\).

若 \(\operatorname{rank} {(A)} < n - 1\),
则 \(A\) 的 \(n-1\)~级子阵 \(A(j|i)\) 的行列式为 \(0\).
则 \(\operatorname{adj} {(A)} = 0\).
则 \(\operatorname{rank} {(\operatorname{adj} {(A)})} = 0\).

有挑战的情形或许是当 \(\operatorname{rank} {(A)} = n - 1\) 时.
我们证明,
\(\operatorname{rank} {(\operatorname{adj} {(A)})} = 1\).

\begin{proof}[证~1]
    我们说, 若 \(\det {(A)} = 0\),
    则 \(\operatorname{adj} {(A)}\)
    没有行列式非零的 \(2\)~级子阵
    (见 ``古伴的性质 (2)'').
    另一方面, 既然 \(\operatorname{rank} {(A)} = n - 1\),
    则 \(A\) 的某 \(n-1\)~级子阵的行列式
    \(\det {(A(v|u))} \neq 0\).
    则 \([\operatorname{adj} {(A)}]_{u,v}
    = (-1)^{v+u} \det {(A(v|u))} \neq 0\).
    则 \(\operatorname{adj} {(A)}\)
    有行列式非零的 \(1\)~级子阵.
    则 \(\operatorname{rank} {(\operatorname{adj} {(A)})} = 1\).
\end{proof}

注意, 倍加不改变秩.
于是, 我们也可考虑用倍加.

\begin{proof}[证~2]
    我们知道, 存在若干个形如
    \(E(n; p, q; s)\)
    的阵
    \(E_1\), \(F_1\), \(E_2\), \(F_2\),
    \(\dots\), \(E_u\), \(F_u\),
    与若干个形如 \(M(n; q'; s')\) 的阵
    \(M_1\), \(M_2\), \(\dots\), \(M_n\),
    使
    \begin{align*}
        A
        = {} &
        F_1 F_2 \dots F_u
        M_1 M_2 \dots M_n
        E_u \dots E_2 E_1
        \\
        = {} &
        F_1 F_2 \dots F_u
        (M_1 M_2 \dots M_n)
        E_u \dots E_2 E_1.
    \end{align*}
    则
    \begin{align*}
        \operatorname{adj} {(A)}
        = {} &
        \operatorname{adj} {(E_1)}
        \dots
        \operatorname{adj} {(E_u)}
        \operatorname{adj} {(M_n)}
        \dots
        \operatorname{adj} {(M_1)}
        \operatorname{adj} {(F_u)}
        \dots
        \operatorname{adj} {(F_1)}
        \\
        = {} &
        \operatorname{adj} {(E_1)}
        \dots
        \operatorname{adj} {(E_u)}
        (\operatorname{adj} {(M_n)}
        \dots
        \operatorname{adj} {(M_1)})
        \operatorname{adj} {(F_u)}
        \dots
        \operatorname{adj} {(F_1)}
        \\
        = {} &
        \operatorname{adj} {(E_1)}
        \dots
        \operatorname{adj} {(E_u)}
        \operatorname{adj} {(M_1 \dots M_n)}
        \operatorname{adj} {(F_u)}
        \dots
        \operatorname{adj} {(F_1)}.
    \end{align*}
    记 \(D = M_1 M_2 \dots M_n\).
    则当 \(i \neq j\) 时, \([D]_{i,j} = 0\).
    则当 \(i \neq j\) 时,
    \([\operatorname{adj} {(D)}]_{i,j} = 0\).
    注意, \(\operatorname{rank} {(A)}
    = \operatorname{rank} {(D)}\).
    注意, \(\operatorname{adj} {(E(n; p, q; s))}
    = E(n; p, q; -s)\),
    故 \(\operatorname{rank} {(\operatorname{adj} {(A)})}
    = \operatorname{rank} {(\operatorname{adj} {(D)})}\).

    设 \(\operatorname{rank} {(A)} = n - 1\).
    则 \(\operatorname{rank} {(D)} = n - 1\).
    设在
    \([D]_{1,1}\), \([D]_{2,2}\), \(\dots\), \([D]_{n,n}\),
    恰有 \(r\)~个数不是 \(0\).
    则 \(\operatorname{rank} {(D)} = r\).
    则在
    \([D]_{1,1}\), \([D]_{2,2}\), \(\dots\), \([D]_{n,n}\),
    恰有 \(n-1\)~个数不是 \(0\).
    设 \([D]_{k,k} = 0\).
    则当 \(i \neq k\) 时, \([D]_{i,i} \neq 0\).
    则在
    \([\operatorname{adj} {(D)}]_{1,1}\),
    \([\operatorname{adj} {(D)}]_{2,2}\),
    \(\dots\),
    \([\operatorname{adj} {(D)}]_{n,n}\),
    有一个数不是 \(0\), 因为
    \begin{align*}
        [\operatorname{adj} {(D)}]_{k,k}
        = (-1)^{k+k} \det {(D(k|k))}
        = \prod_{\substack{1 \leq \ell \leq n \\\ell \neq k}}
        {[D]_{\ell,\ell}}
        \neq 0.
    \end{align*}
    并且, 当 \(i \neq k\) 时,
    \begin{align*}
        [\operatorname{adj} {(D)}]_{i,i}
        = (-1)^{i+i} \det {(D(i|i))}
        = \prod_{\substack{1 \leq \ell \leq n \\\ell \neq i}}
        {[D]_{\ell,\ell}}
            = [D]_{k,k}
        \prod_{\substack{1 \leq \ell \leq n   \\\ell \neq i, k}}
        {[D]_{\ell,\ell}}
        = 0.
    \end{align*}
    则 \(\operatorname{rank} {(\operatorname{adj} {(D)})} = 1\).
    则 \(\operatorname{rank} {(\operatorname{adj} {(A)})} = 1\).
\end{proof}

\section{REF 阵、线性方程组、秩}

本节, 我们讨论 REF 阵,
并用它研究线性方程组与秩.

我们知道, 最大的整数不存在:
对任何整数 \(x\),
我们能作大于 \(x\) 的整数,
如 \(x + 1\).
但我们可作 ``Tiānyī 记号'',
且考虑它与数的关系:

\begin{definition}
    我们定义 \emph{Tiānyī 记号}为 \(\infty\).
    Tiānyī 记号 \(\infty\) 可被叫作 Tiānyī.

    我们规定, 对任何整数 \(x\),
    Tiānyī ``不等于'' \(x\),
    且 \(x\) ``不等于'' Tiānyī:
    \begin{align*}
        \infty \neq x
        \quad \text{且} \quad
        x \neq \infty.
    \end{align*}

    我们规定, 对任何整数 \(x\),
    Tiānyī ``大于'' \(x\),
    且 \(x\) ``小于'' Tiānyī:
    \begin{align*}
        \infty > x
        \quad \text{且} \quad
        x < \infty.
    \end{align*}

    我们规定, Tiānyī ``等于'' Tiānyī:
    \begin{align*}
        \infty = \infty.
    \end{align*}
\end{definition}

Tiānyī 记号来自中国 Vocaloid \pinjino{bokaloyido}
\href{https://zh.moegirl.org.cn/洛天依}{Luò Tiānyī}.

注意, \(\geq\), \(\leq\) 分别是
``大于或等于'', ``小于或等于''.
于是, 对任何整数 \(x\),
我们当然有 \(\infty \geq x\)
与 \(x \leq \infty\).

现在, 我们可方便地定义 REF 阵.

\begin{definition}
    设 \(A\) 是 \(m \times n\)~阵.

    设 \(i\) 是不超过 \(m\) 的正整数.
    若 \(A\) 的行~\(i\) 的元全为 \(0\),
    我们说, \(A\) 的行~\(i\) 的 \emph{pivot} \pinjino{pivite}
    是 Tiānyī:
    \begin{align*}
        \operatorname{pivot} {(A; i)} = \infty.
    \end{align*}
    若 \(A\) 的行~\(i\) 有不为 \(0\) 的元,
    则一定存在不超过 \(n\) 的正整数 \(k_i\),
    使 \([A]_{i,k_i} \neq 0\),
    且当 \(j < k_i\) 时, \([A]_{i,j} = 0\).
    我们说, \(A\) 的行~\(i\) 的 \emph{pivot} 是 \(k_i\):
    \begin{align*}
        \operatorname{pivot} {(A; i)} = k_i.
    \end{align*}

    注意, \(m \times n\)~阵 \(A\) 有 \(m\) pivots.

    设存在不高于 \(m\) 的非负整数 \(r\), 使:

    (1)
    对任何不高于 \(r\) 的正整数 \(i\),
    \begin{align*}
        \operatorname{pivot} {(A; i)} < \infty;
    \end{align*}

    (2)
    对任何低于 \(r\) 的正整数 \(i\),
    \begin{align*}
        \operatorname{pivot} {(A; i)}
        < \operatorname{pivot} {(A; i + 1)};
    \end{align*}

    (3)
    对任何高于 \(r\) 且不高于 \(m\) 的整数 \(i\),
    \begin{align*}
        \operatorname{pivot} {(A; i)} = \infty.
    \end{align*}

    我们说, \(A\) 是 \emph{REF}
    (\angla{row-echelon form}) \emph{阵}.
\end{definition}

注意, 当 \(r = 0\) 时,
不存在不高于 \(r\) 的正整数
(且不存在低于 \(r\) 的正整数),
故条件~(1) 与 (2) 成立.
则 \(A\) 的所有的 pivots 是 Tiānyī,
故 \(A = 0\).

注意, 当 \(r = 1\) 时,
不存在低于 \(r\) 的正整数,
故条件~(2) 成立.

注意, \(1 \times n\)~阵是 REF 阵.

注意, 当 \(r = m\) 时,
不存在整数, 其高于 \(r\) 且不高于 \(m\),
故条件~(3) 成立.

\begin{example}
    零阵是 REF 阵 (\(r = 0\)).
    \(1 \times n\)~阵也是 REF 阵.

    下面的阵是 REF 阵, 其中,
    我们标记了对应非 Tiānyī 的 pivots 的元:
    \begin{align*}
        \begin{bmatrix}
            \boxed{6} & 8         & 15         \\
            0         & \boxed{4} & 15         \\
            0         & 0         & \boxed{15} \\
        \end{bmatrix},
        \quad
        \begin{bmatrix}
            \boxed{37} & 0          & 73 & 0          & 1  \\
            0          & \boxed{39} & 54 & 93         & 0  \\
            0          & 0          & 0  & \boxed{86} & 20 \\
            0          & 0          & 0  & 0          & 0  \\
        \end{bmatrix},
        \quad
        \begin{bmatrix}
            0 & \boxed{6} & 2 & 8         \\
            0 & 0         & 0 & \boxed{3} \\
            0 & 0         & 0 & 0         \\
            0 & 0         & 0 & 0         \\
            0 & 0         & 0 & 0         \\
            0 & 0         & 0 & 0         \\
        \end{bmatrix}.
    \end{align*}

    下面的阵不是 REF 阵, 其中,
    我们标记了对应非 Tiānyī 的 pivots 的元:
    \begin{align*}
        \begin{bmatrix}
            \boxed{1} \\
            0         \\
            \boxed{1} \\
        \end{bmatrix},
        \quad
        \begin{bmatrix}
            \boxed{37} & 73         & 0  & 0  \\
            0          & \boxed{39} & 0  & 93 \\
            0          & \boxed{86} & 20 & 0  \\
            0          & 0          & 0  & 0  \\
        \end{bmatrix},
        \quad
        \begin{bmatrix}
            0 & 0         & 0         & \boxed{6} \\
            0 & 0         & \boxed{2} & 0         \\
            0 & \boxed{8} & 0         & 0         \\
        \end{bmatrix}.
    \end{align*}
\end{example}

\begin{theorem}
    设 \(A\) 是 \(m \times n\)~阵.
    则我们可用行的倍加,
    变 \(A\) 为某 \(m \times n\) REF 阵 \(B\).
\end{theorem}

\begin{proof}
    我们用数学归纳法证明此事.
    具体地, 设 \(P(m)\) 为命题
    \begin{quotation}
        对任何正整数 \(m\),
        对任何 \(m \times n\)~阵 \(A\),
        用若干次行的倍加,
        我们可变 \(A\) 为 \(m \times n\)~阵 \(B\),
        使 \(B\) 是 REF 阵.
    \end{quotation}
    则, 我们的目标是,
    对任何正整数 \(m\), \(P(m)\) 是对的.

    \(P(1)\) 是对的.

    假定 \(P(m-1)\) 是对的.
    我们由此证 \(P(m)\) 也是对的.

    若 \(A = 0\), 则 \(A\) 已是 REF 阵.

    设 \(A \neq 0\).
    则 \(A\) 有最小的 pivot,
    且它小于 Tiānyī;
    具体地, 存在不超过 \(m\) 的正整数 \(\ell\),
    使对任何不超过 \(m\) 的正整数 \(i\),
    \(\operatorname{pivot} {(A; \ell)}
    \leq \operatorname{pivot} {(A; i)}\),
    且
    \(\operatorname{pivot} {(A; \ell)} < \infty\).

    若 \(\ell \neq 1\),
    我们可用三次行的倍加,
    以 ``几乎'' 交换行~\(1\) 与行~\(\ell\),
    得 \(m \times n\)~阵 \(A_1\):
    \begin{align*}
        (x, y)
        \to {} & (x + 1y, y) = (x + y, y)                \\
        \to {} & (x + y, y + (-1) (x + y)) = (x + y, -x) \\
        \to {} & (x + y + 1(-x), -x) = (y, -x).
    \end{align*}
    具体地,
    \begin{align*}
        [A_1]_{i,j} =
        \begin{cases}
            [A]_{\ell,j}, & i = 1;     \\
            -[A]_{1,j},   & i = \ell;  \\
            [A]_{i,j},    & \text{其他}.
        \end{cases}
    \end{align*}
    则
    \(\operatorname{pivot} {(A_1; 1)}
    = \operatorname{pivot} {(A; \ell)}\),
    \(\operatorname{pivot} {(A_1; \ell)}
    = \operatorname{pivot} {(A; 1)}\)
    (以非零的数乘一行不改变它的 pivot),
    且
    \(\operatorname{pivot} {(A_1; i)}
    = \operatorname{pivot} {(A; i)}\)
    (若 \(i \neq \ell\) 或 \(1\)).

    若 \(\ell = 1\), 则我们不变, 取 \(A_1\) 为 \(A\).

    综上, 我们可用行的倍加,
    变 \(A\) 为 \(A_1\),
    使 \(A_1\) 的行~\(1\) 有最小的 pivot,
    且它小于 Tiānyī.

    记 \(k_1 = \operatorname{pivot} {(A_1; 1)}\).
    注意, 既然 \(A_1\) 的行~\(1\) 有最小的 pivot,
    且它小于 Tiānyī,
    那么, 对 \(j < k_1\),
    必 \([A_1]_{i,j} = 0\).
    注意, \([A_1]_{1,k_1} \neq 0\).
    那么, 对 \(A_1\),
    我们%
    加行~\(1\) 的 \(-[A_1]_{2,k_1}/[A_1]_{1,k_1}\) 倍于行~\(2\),
    加行~\(1\) 的 \(-[A_1]_{3,k_1}/[A_1]_{1,k_1}\) 倍于行~\(3\),
    ……
    加行~\(1\) 的 \(-[A_1]_{m,k_1}/[A_1]_{1,k_1}\) 倍于行~\(m\),
    得阵 \(C\).
    则对 \(j < k_1\), 必 \([C]_{i,j} = 0\),
    且 \([C]_{1,k_1} \neq 0\),
    且 \([C]_{i,k_1} = 0\), 若 \(i > 1\).
    则 \(C\) 的行~\(1\) 有最小的 pivot \(k_1\),
    且它小于 Tiānyī.

    若 \(k_1 = n\), 则 \(C\) 已是 REF 阵.

    设 \(k_1 < n\).
    考虑 \((m - 1) \times (n - k_1)\)~阵
    \(C(1|{1,2,\dots,k_1})\).
    由假定, 用若干次行的倍加,
    我们可变 \(C(1|{1,2,\dots,k_1})\) 为
    \((m - 1) \times (n - k_1)\)~阵 \(T\),
    使 \(T\) 是 REF 阵.

    注意, 既然当 \(1 < i\), 且 \(j \leq k_1\) 时,
    \([C]_{i,j} = 0\),
    那么, 无论如何对 \(C\) 的不是行~\(1\) 的行作倍加,
    得到的阵的 \((i, j)\)-元是 \(0\)
    (\(1 < i\), 且 \(j \leq k_1\)).
    那么, 作若干次行的倍加后,
    我们可变 \(C\) 为 \(m \times n\)~阵 \(B\), 使
    \begin{align*}
        [B]_{i,j} =
        \begin{cases}
            [T]_{i-1,j-k_1},
             & \text{\(i > 1\), 且 \(j > k_1\)}; \\
            [C]_{i,j},
             & \text{其他}.
        \end{cases}
    \end{align*}
    注意, 对 \(i > 1\),
    \begin{align*}
        \operatorname{pivot} {(B; i)} =
        \begin{cases}
            \operatorname{pivot} {(T; i-1)} + k_1,
             & \operatorname{pivot} {(T; i-1)} < \infty;
            \\
            \infty,
             & \operatorname{pivot} {(T; i-1)} = \infty.
        \end{cases}
    \end{align*}
    由此可见, \(B\) 是 REF 阵.

    所以, \(P(m)\) 是对的.
    由数学归纳法, 待证命题成立.
\end{proof}

\begin{example}
    记
    \begin{align*}
        A_1 =
        \begin{bmatrix}
            1 \\
            0 \\
            1 \\
        \end{bmatrix},
        \quad
        A_2 =
        \begin{bmatrix}
            37 & 73 & 0  & 0  \\
            0  & 39 & 0  & 93 \\
            0  & 86 & 20 & 0  \\
            0  & 0  & 0  & 0  \\
        \end{bmatrix},
        \quad
        A_3 =
        \begin{bmatrix}
            0 & 0 & 0 & 6 \\
            0 & 0 & 2 & 0 \\
            0 & 8 & 0 & 0 \\
        \end{bmatrix}.
    \end{align*}

    对 \(A_1\),
    我们加行~\(1\) 的 \(-1\)~倍于行~\(3\), 得
    \begin{align*}
        \begin{bmatrix}
            1 \\
            0 \\
            0 \\
        \end{bmatrix}.
    \end{align*}
    这是 REF 阵.

    对 \(A_2\),
    我们加行~\(2\) 的 \(-86/39\)~倍于行~\(3\), 得
    \begin{align*}
        \begin{bmatrix}
            37 & 73 & 0  & 0          \\
            0  & 39 & 0  & 93         \\
            0  & 0  & 20 & -2\,666/13 \\
            0  & 0  & 0  & 0          \\
        \end{bmatrix}.
    \end{align*}
    这是 REF 阵.

    对 \(A_3\),
    我们用三次行的倍加,
    ``交换'' 行~\(1\) 与行~\(3\), 得
    \begin{align*}
        \begin{bmatrix}
            0 & 8 & 0 & 0  \\
            0 & 0 & 2 & 0  \\
            0 & 0 & 0 & -6 \\
        \end{bmatrix}.
    \end{align*}
    这是 REF 阵.
\end{example}

\begin{example}
    设
    \begin{align*}
        G =
        \begin{bmatrix}
            0 & 1 & 1 & 1 & 1 & 17 \\
            1 & 0 & 1 & 1 & 1 & 15 \\
            1 & 1 & 0 & 1 & 1 & 14 \\
            1 & 1 & 1 & 0 & 1 & 12 \\
            1 & 1 & 1 & 1 & 0 & 10 \\
        \end{bmatrix}.
    \end{align*}
    我们用行的倍加变它为 REF 阵.

    注意, 当 \(j \leq 5\) 时,
    \(G\) 的列~\(j\) 的 \(5\) 个元的和为 \(4\).
    于是, 我们考虑分别地加行
    \(2\), \(3\), \(4\), \(5\)
    于行~\(1\), 得
    \begin{align*}
        \begin{bmatrix}
            4 & 4 & 4 & 4 & 4 & 68 \\
            1 & 0 & 1 & 1 & 1 & 15 \\
            1 & 1 & 0 & 1 & 1 & 14 \\
            1 & 1 & 1 & 0 & 1 & 12 \\
            1 & 1 & 1 & 1 & 0 & 10 \\
        \end{bmatrix}.
    \end{align*}
    然后, 我们分别地加行~\(1\) 的 \(-1/4\)~倍于行
    \(2\), \(3\), \(4\), \(5\), 得
    \begin{align*}
        \begin{bmatrix}
            4 & 4  & 4  & 4  & 4  & 68 \\
            0 & -1 & 0  & 0  & 0  & -2 \\
            0 & 0  & -1 & 0  & 0  & -3 \\
            0 & 0  & 0  & -1 & 0  & -5 \\
            0 & 0  & 0  & 0  & -1 & -7 \\
        \end{bmatrix}.
    \end{align*}
    这是 REF 阵.
    当然, 若我们分别地加行
    \(2\), \(3\), \(4\), \(5\)
    的 \(4\)~倍于行~\(1\), 则我们有
    \begin{align*}
        \begin{bmatrix}
            4 & 0  & 0  & 0  & 0  & 0  \\
            0 & -1 & 0  & 0  & 0  & -2 \\
            0 & 0  & -1 & 0  & 0  & -3 \\
            0 & 0  & 0  & -1 & 0  & -5 \\
            0 & 0  & 0  & 0  & -1 & -7 \\
        \end{bmatrix}.
    \end{align*}
    这仍是 REF 阵.
\end{example}

现在, 我们讨论线性方程组.

回想,
线性方程组是形如
\begin{equation*}
    \left\{
    \begin{aligned}
        a_{1,1} x_1 + a_{1,2} x_2 + \dots + a_{1,n} x_n & = b_1, \\
        a_{2,1} x_1 + a_{2,2} x_2 + \dots + a_{2,n} x_n & = b_2, \\
                                                        & \dots, \\
        a_{m,1} x_1 + a_{m,2} x_2 + \dots + a_{m,n} x_n & = b_m
    \end{aligned}
    \right.
\end{equation*}
的方程组,
其中,
\(a_{1,1}\), \(a_{1,2}\), \(\dots\), \(a_{1,n}\),
\(a_{2,1}\), \(a_{2,2}\), \(\dots\), \(a_{2,n}\),
\(\dots\),
\(a_{m,1}\), \(a_{m,2}\), \(\dots\), \(a_{m,n}\),
\(b_1\), \(b_2\), \(\dots\), \(b_m\)
是事先指定的 \(mn + m\) 个数,
且 \(x_1\), \(x_2\), \(\dots\), \(x_n\)
是未知数.
若数 \(c_1\), \(c_2\), \(\dots\), \(c_n\) 适合,
\begin{align*}
    a_{1,1} c_1 + a_{1,2} c_2 + \dots + a_{1,n} c_n & = b_1, \\
    a_{2,1} c_1 + a_{2,2} c_2 + \dots + a_{2,n} c_n & = b_2, \\
                                                    & \dots, \\
    a_{m,1} c_1 + a_{m,2} c_2 + \dots + a_{m,n} c_n & = b_m,
\end{align*}
我们说,
\(x_1 = c_1\),
\(x_2 = c_2\),
\(\dots\),
\(x_n = c_n\)
是此方程组的一个解.

我们可用若干次 ``加一个方程的倍于另一个方程'' 的行为,
变一个线性方程组为简单的线性方程组,
且它们有相同的解.

\begin{theorem}
    考虑线性方程组
    \begin{equation}
        \left\{
        \begin{aligned}
            a_{1,1} x_1 + a_{1,2} x_2 + \dots + a_{1,n} x_n & = b_1, \\
            a_{2,1} x_1 + a_{2,2} x_2 + \dots + a_{2,n} x_n & = b_2, \\
                                                            & \dots, \\
            a_{m,1} x_1 + a_{m,2} x_2 + \dots + a_{m,n} x_n & = b_m.
        \end{aligned}
        \right.
        \label{eq:sole-1}
    \end{equation}
    设 \(p\), \(q\) 是不超过 \(m\) 的正整数,
    且互不相同.
    设 \(s\) 是数.
    加方程~\(p\) 的 \(s\)~倍于方程~\(q\),
    且不改变其他的方程,
    得线性方程组
    \begin{equation}
        \left\{
        \begin{aligned}
            a'_{1,1} x_1 + a'_{1,2} x_2 + \dots + a'_{1,n} x_n & = b'_1, \\
            a'_{2,1} x_1 + a'_{2,2} x_2 + \dots + a'_{2,n} x_n & = b'_2, \\
                                                               & \dots,  \\
            a'_{m,1} x_1 + a'_{m,2} x_2 + \dots + a'_{m,n} x_n & = b'_m,
        \end{aligned}
        \right.
        \label{eq:sole-2}
    \end{equation}
    其中,
    \begin{align*}
        a'_{i,j} =
        \begin{cases}
            a_{i,j},             & i \neq q; \\
            a_{q,j} + s a_{p,j}, & i = q,
        \end{cases}
    \end{align*}
    且
    \begin{align*}
        b'_i =
        \begin{cases}
            b_i,         & i \neq q; \\
            b_q + s b_p, & i = q.
        \end{cases}
    \end{align*}
    则方程组~\eqref{eq:sole-1} 的解%
    都是方程组~\eqref{eq:sole-2} 的解,
    且方程组~\eqref{eq:sole-2} 的解%
    都是方程组~\eqref{eq:sole-1} 的解.
\end{theorem}

\begin{proof}
    设数 \(c_1\), \(c_2\), \(\dots\), \(c_n\) 适合
    \begin{align*}
        a_{1,1} c_1 + a_{1,2} c_2 + \dots + a_{1,n} c_n & = b_1, \\
        a_{2,1} c_1 + a_{2,2} c_2 + \dots + a_{2,n} c_n & = b_2, \\
                                                        & \dots, \\
        a_{m,1} c_1 + a_{m,2} c_2 + \dots + a_{m,n} c_n & = b_m.
    \end{align*}
    那么, 若 \(i \neq q\), 则
    \begin{align*}
             &
        a'_{i,1} c_1 + a'_{i,2} c_2 + \dots + a'_{i,n} c_n
        \\
        = {} &
        a_{i,1} c_1 + a_{i,2} c_2 + \dots + a_{i,n} c_n
        \\
        = {} &
        b_i
        \\
        = {} &
        b'_i;
    \end{align*}
    若 \(i = q\), 则
    \begin{align*}
             &
        a'_{i,1} c_1 + a'_{i,2} c_2 + \dots + a'_{i,n} c_n
        \\
        = {} &
        (a_{q,1} + s a_{p,1}) c_1
        + (a_{q,2} + s a_{p,2}) c_2
        + \dots
        + (a_{q,n} + s a_{p,n}) c_n
        \\
        = {} &
        (a_{q,1} c_1 + a_{q,2} c_2 + \dots + a_{q,n} c_n)
        + s (a_{p,1} c_1 + a_{p,2} c_2 + \dots + a_{p,n} c_n)
        \\
        = {} &
        b_q + s b_p
        \\
        = {} &
        b'_i.
    \end{align*}
    则方程组~\eqref{eq:sole-1} 的解%
    都是方程组~\eqref{eq:sole-2} 的解.

    对方程组~\eqref{eq:sole-2},
    我们加方程~\(p\) 的 \(-s\) 倍于方程~\(q\),
    且不改变其他的方程,
    得方程组
    \begin{equation}
        \left\{
        \begin{aligned}
            a''_{1,1} x_1 + a''_{1,2} x_2 + \dots + a''_{1,n} x_n & = b''_1, \\
            a''_{2,1} x_1 + a''_{2,2} x_2 + \dots + a''_{2,n} x_n & = b''_2, \\
                                                                  & \dots,   \\
            a''_{m,1} x_1 + a''_{m,2} x_2 + \dots + a''_{m,n} x_n & = b''_m,
        \end{aligned}
        \right.
        \label{eq:sole-3}
    \end{equation}
    其中,
    \begin{align*}
        a''_{i,j} =
        \begin{cases}
            a'_{i,j},                 & i \neq q; \\
            a'_{q,j} + (-s) a'_{p,j}, & i = q,
        \end{cases}
    \end{align*}
    且
    \begin{align*}
        b''_i =
        \begin{cases}
            b'_i,             & i \neq q; \\
            b'_q + (-s) b'_p, & i = q.
        \end{cases}
    \end{align*}
    则当 \(i \neq q\) 时,
    \begin{align*}
        a''_{i,j} & = a'_{i,j} = a_{i,j}, \\
        b''_i     & = b'_i = b_i,
    \end{align*}
    且当 \(i = q\) 时,
    \begin{align*}
        a''_{i,j}
         & = a'_{q,j} - s a'_{p,j}
        = a_{q,j} + s a_{p,j} - s a_{p,j}
        = a_{i,j},                 \\
        b''_i
         & = b'_q - s b'_p
        = b_q + s b_p - s b_p
        = b_i.
    \end{align*}
    故方程组~\eqref{eq:sole-3} 就是方程组~\eqref{eq:sole-1}.
    那么, 由前面的讨论,
    方程组~\eqref{eq:sole-2} 的解%
    都是方程组~\eqref{eq:sole-3} 的解,
    即方程组~\eqref{eq:sole-1} 的解.
\end{proof}

考虑线性方程组
\begin{equation}
    \left\{
    \begin{aligned}
        a_{1,1} x_1 + a_{1,2} x_2 + \dots + a_{1,n} x_n & = b_1, \\
        a_{2,1} x_1 + a_{2,2} x_2 + \dots + a_{2,n} x_n & = b_2, \\
                                                        & \dots, \\
        a_{m,1} x_1 + a_{m,2} x_2 + \dots + a_{m,n} x_n & = b_m.
    \end{aligned}
    \right.
    \label{eq:sole-4}
\end{equation}
我们说, 它跟 \(m \times (n + 1)\)~阵
\begin{equation}
    G =
    \begin{bmatrix}
        a_{1,1} & a_{1,2} & \cdots & a_{1,n} & b_1    \\
        a_{2,1} & a_{2,2} & \cdots & a_{2,n} & b_2    \\
        \vdots  & \vdots  & {}     & \vdots  & \vdots \\
        a_{m,1} & a_{m,2} & \cdots & a_{m,n} & b_m    \\
    \end{bmatrix}
    \label{eq:sole-4-mat}
\end{equation}
对应:
给定线性方程组~\eqref{eq:sole-4},
我们能以 \(m \times (n + 1)\)~阵表示它,
且反过来, 给定 \(m \times (n + 1)\)~阵 \eqref{eq:sole-4-mat},
我们能还原它为线性方程组~\eqref{eq:sole-4}.
并且, 对线性方程组~\eqref{eq:sole-4},
加方程~\(p\) 的 \(s\)~倍于方程~\(q\) (\(q \neq p\)),
且不改变其他的方程,
相当于对阵~\eqref{eq:sole-4-mat},
加行~\(p\) 的 \(s\)~倍于行~\(q\)  (\(q \neq p\)),
且不改变其他的行.

我们知道, 我们可用若干次行的倍加,
变 \(G\) 为 REF 阵 \(G'\).
则存在非负整数 \(u\), 使
\begin{align*}
    \operatorname{pivot} {(G'; 1)}
    < \operatorname{pivot} {(G'; 2)}
    < \dots
    < \operatorname{pivot} {(G'; u-1)}
    < \operatorname{pivot} {(G'; u)}
    < \infty,
\end{align*}
且对 \(i > u\),
\(G'\) 的行~\(i\) 是 \(0\).
记 \(\operatorname{pivot} {(G'; i)} = p_i\)
(\(i = 1\), \(2\), \(\dots\), \(u\)).

若 \(p_u = n + 1\),
则 \(G'\) 形如
\begin{align*}
    \begin{bmatrix}
        {}         & [G']_{1,p_1} & \cdots & [G']_{1,p_2} & \cdots     & [G']_{1,p_{u-1}}   & \cdots & [G']_{1,n+1}   \\
        {}         & 0            & \cdots & [G']_{2,p_2} & \cdots     & [G']_{2,p_{u-1}}   & \cdots & [G']_{2,n+1}   \\
        \mathsf{O} & \vdots       & {}     & \vdots       & {}         & \vdots             & \cdots & \vdots         \\
        {}         & 0            & \cdots & 0            & \cdots     & [G']_{u-1,p_{u-1}} & \cdots & [G']_{u-1,n+1} \\
        {}         & 0            & \cdots & 0            & \cdots     & 0                  & \cdots & [G']_{u,n+1}   \\
        {}         & {}           & {}     & {}           & \mathsf{O} & {}                 & {}     & {}             \\
    \end{bmatrix},
\end{align*}
其中, \(\mathsf{O}\) 表示可能出现的 \(0\), 下同.
记 \(r = u - 1\),
且 \(k_i = p_i\)
(\(i = 1\), \(2\), \(\dots\), \(u - 1\)).

若 \(p_u \leq n\),
则 \(G'\) 形如
\begin{align*}
    \begin{bmatrix}
        {}         & [G']_{1,p_1} & \cdots & [G']_{1,p_2} & \cdots     & [G']_{1,p_u} & \cdots & [G']_{1,n+1} \\
        {}         & 0            & \cdots & [G']_{2,p_2} & \cdots     & [G']_{2,p_u} & \cdots & [G']_{2,n+1} \\
        \mathsf{O} & \vdots       & {}     & \vdots       & {}         & \vdots       & \cdots & \vdots       \\
        {}         & 0            & \cdots & 0            & \cdots     & [G']_{u,p_u} & \cdots & [G']_{u,n+1} \\
        {}         & {}           & {}     & {}           & \mathsf{O} & {}           & {}     & {}           \\
    \end{bmatrix}.
\end{align*}
记 \(r = u\),
且 \(k_i = p_i\)
(\(i = 1\), \(2\), \(\dots\), \(u\)).

于是, 统一地, \(G'\) 形如
\begin{align*}
    \begin{bmatrix}
        {}         & [G']_{1,k_1} & \cdots & [G']_{1,k_2} & \cdots     & [G']_{1,k_r} & \cdots & [G']_{1,n+1}   \\
        {}         & 0            & \cdots & [G']_{2,k_2} & \cdots     & [G']_{2,k_r} & \cdots & [G']_{2,n+1}   \\
        \mathsf{O} & \vdots       & {}     & \vdots       & {}         & \vdots       & \cdots & \vdots         \\
        {}         & 0            & \cdots & 0            & \cdots     & [G']_{r,k_r} & \cdots & [G']_{r,n+1}   \\
        {}         & 0            & \cdots & 0            & \cdots     & 0            & \cdots & [G']_{r+1,n+1} \\
        {}         & {}           & {}     & {}           & \mathsf{O} & {}           & {}     & {}             \\
    \end{bmatrix},
\end{align*}
其中,
\(1 \leq k_1 < k_2 < \dots < k_r \leq n\),
且 \([G']_{i,k_i} \neq 0\)
(\(i = 1\), \(2\), \(\dots\), \(r\)),
且当 \(i > r + 1\) 时, \([G']_{i,n+1} = 0\),
且当 \(i > r\), 且 \(j \leq n\) 时,
\([G']_{i,j} = 0\).
(注意, 若 \(r = m\), 则 \([G']_{r+1,n+1}\) 不出现.)

% 我们知道, 我们可用行的倍加, 变一个阵为一个 REF 阵.
% 设我们用行的倍加, 变阵~\eqref{eq:sole-4-mat}
% 为 REF 阵
% \begin{align*}
%     G' =
%     \begin{bmatrix}
%         {}         & a'_{1,k_1} & \cdots & a'_{1,k_2} & \cdots     & a'_{1,k_r} & \cdots & b'_1     \\
%         {}         & 0          & \cdots & a'_{2,k_2} & \cdots     & a'_{2,k_r} & \cdots & b'_2     \\
%         \mathsf{O} & \vdots     & {}     & \vdots     & {}         & \vdots     & \cdots & \vdots   \\
%         {}         & 0          & \cdots & 0          & \cdots     & a'_{r,k_r} & \cdots & b'_r     \\
%         {}         & 0          & \cdots & 0          & \cdots     & 0          & \cdots & b'_{r+1} \\
%         {}         & {}         & {}     & {}         & \mathsf{O} & {}         & {}     & {}       \\
%     \end{bmatrix},
% \end{align*}
% 其中,
% \(\mathsf{O}\) 表示可能出现的 \(0\);
% 具体地, 存在不高于 \(m\) 的非负整数 \(r\), 使
% \begin{align*}
%     \operatorname{pivot} {(G'; i)} = k_i \leq n,
%       & \quad
%     \text{\(i = 1\), \(2\), \(\dots\), \(r\)};
%     \\
%     \operatorname{pivot} {(G'; i)}
%     < \operatorname{pivot} {(G'; i + 1)},
%       & \quad
%     \text{\(i = 1\), \(\dots\), \(r-1\)};
%     \\
%     \operatorname{pivot} {(G'; r+1)}
%     = &
%     \begin{cases}
%         n + 1,  & b'_{r+1} \neq 0; \\
%         \infty, & b'_{r+1} = 0;
%     \end{cases}
%     \\
%     \operatorname{pivot} {(G'; i)} = \infty,
%       & \quad
%     \text{\(i > r + 1\)}.
% \end{align*}
% (注意, 若 \(r = m\),
% 则行~\(r + 1\), \(r + 2\), \(\dots\) 当然不出现.)

为方便, 记 \([G']_{i,n+1} = b'_i\).
则当 \(i > r + 1\) 时, \(b'_i = 0\).
记 \([G']_{i,j} = a'_{i,j}\)
(若 \(j \leq n\)).
则当 \(i > r\), 且 \(j \leq n\) 时,
\(a'_{i,j} = 0\).
则 \(G'\) 形如
\begin{align*}
    \begin{bmatrix}
        {}         & a'_{1,k_1} & \cdots & a'_{1,k_2} & \cdots     & a'_{1,k_r} & \cdots & b'_1     \\
        {}         & 0          & \cdots & a'_{2,k_2} & \cdots     & a'_{2,k_r} & \cdots & b'_2     \\
        \mathsf{O} & \vdots     & {}     & \vdots     & {}         & \vdots     & \cdots & \vdots   \\
        {}         & 0          & \cdots & 0          & \cdots     & a'_{r,k_r} & \cdots & b'_r     \\
        {}         & 0          & \cdots & 0          & \cdots     & 0          & \cdots & b'_{r+1} \\
        {}         & {}         & {}     & {}         & \mathsf{O} & {}         & {}     & {}       \\
    \end{bmatrix}.
\end{align*}
则方程组~\eqref{eq:sole-4} 被相应地变为
\begin{equation}
    \left\{
    \begin{aligned}
        a'_{1,k_1} x_{k_1} + \sum_{k_1 < j \leq n} {a'_{1,j} x_{j}} & = b'_1,     \\
        a'_{2,k_2} x_{k_2} + \sum_{k_2 < j \leq n} {a'_{2,j} x_{j}} & = b'_2,     \\
                                                                    & \dots,      \\
        a'_{r,k_r} x_{k_r} + \sum_{k_r < j \leq n} {a'_{r,j} x_{j}} & = b'_r,     \\
        0                                                           & = b'_{r+1}, \\
        0                                                           & = 0,        \\
                                                                    & \dots,      \\
        0                                                           & = 0,
    \end{aligned}
    \right.
    \label{eq:sole-5}
\end{equation}
其中, \(1 \leq k_1 < k_2 < \dots < k_r \leq n\),
且 \(a'_{i,k_i} \neq 0\)
(\(i = 1\), \(2\), \(\dots\), \(r\)).
(注意, 若 \(r = m\),
则方程 ``\(0 = b'_{r+1}\)'' 当然不出现.)
由前面的讨论, 方程组~\eqref{eq:sole-5}
与方程组~\eqref{eq:sole-4} 有相同的解.

(1)
设 \(b'_{r+1} \neq 0\).
则方程组~\eqref{eq:sole-5} 的方程
``\(0 = b'_{r+1}\)'' 是矛盾.
则方程组~\eqref{eq:sole-5} 无解.
则方程组~\eqref{eq:sole-4} 也无解.
(注意, 此时必 \(r < m\).)

(2)
设 \(b'_{r+1} = 0\) (当 \(r < m\)),
或 ``\(0 = b'_{r+1}\)'' 不出现 (当 \(r = m\)).

若 \(r = n\),
则 \(k_i = i\) (\(i = 1\), \(2\), \(\dots\), \(n\)).
则方程组~\eqref{eq:sole-5} 形如
\begin{equation*}
    \left\{
    \begin{aligned}
        a'_{1,1} x_1 + a'_{1,2} x_2 + \dots + a'_{1,n} x_n & = b'_1, \\
        a'_{2,2} x_2 + \dots + a'_{2,n} x_n                & = b'_2, \\
                                                           & \dots,  \\
        a'_{n,n} x_n                                       & = b'_n,
    \end{aligned}
    \right.
\end{equation*}
其中, \(a'_{i,i} \neq 0\).
这里, 我们未写形如 ``\(0 = 0\)'' 的方程 (若它们出现),
因为它们是恒等式,
且去除它们不影响方程组的解.
由方程~\(n\), 我们有
\begin{align*}
    a'_{n,n} x_n = b'_n,
    \quad
    x_n = \frac{b'_n}{a'_{n,n}};
\end{align*}
代入它到方程~\(n-1\), 我们有
\begin{align*}
    a'_{n-1,n-1} x_{n-1} + a'_{n-1,n} x_{n} = b'_{n-1},
    \quad
    x_{n-1} = \frac{b'_{n-1} - a'_{n-1,n} x_{n}}{a'_{n-1,n-1}};
\end{align*}
……
最后, 我们有
\begin{align*}
    a'_{1,1} x_1 + a'_{1,2} x_2 + \dots + a'_{1,n} x_n = b'_1,
    \quad
    x_1 = \frac{b'_1 - (a'_{1,2} x_2 + \dots + a'_{1,n} x_n)}{a'_{1,1}}.
\end{align*}
则方程组~\eqref{eq:sole-5} 有唯一的解.
则方程组~\eqref{eq:sole-4} 有唯一的解.

若 \(r < n\), 则从 \(1\), \(2\), \(\dots\), \(n\)
去除 \(k_1\), \(\dots\), \(k_r\) 后,
还剩 \(n - r\)~个数.
我们从小到大地叫这 \(n - r\)~个数为
\(\ell_1\), \(\dots\), \(\ell_{n-r}\).
于是, 我们移含
\(x_{\ell_1}\), \(\dots\), \(x_{\ell_{n-r}}\) 的项到右侧,
有
\begin{equation*}
    \left\{
    \begin{aligned}
        a'_{1,k_1} x_{k_1} + a'_{1,k_2} x_{k_2} + \dots + a'_{1,k_r} x_{k_r} & = b'_1 - (a'_{1,\ell_1} x_{\ell_1} + \dots + a'_{1,\ell_{n-r}} x_{\ell_{n-r}}), \\
        a'_{2,k_2} x_{k_2} + \dots + a'_{2,k_r} x_{k_r}                      & = b'_2 - (a'_{2,\ell_1} x_{\ell_1} + \dots + a'_{2,\ell_{n-r}} x_{\ell_{n-r}}), \\
                                                                             & \dots,                                                                          \\
        a'_{r,k_r} x_{k_r}                                                   & = b'_r - (a'_{r,\ell_1} x_{\ell_1} + \dots + a'_{r,\ell_{n-r}} x_{\ell_{n-r}}),
    \end{aligned}
    \right.
\end{equation*}
其中, \(a'_{i,k_i} \neq 0\).
这里, 我们未写形如 ``\(0 = 0\)'' 的方程 (若它们出现).
选取 \(x_{\ell_1}\), \(\dots\), \(x_{\ell_{n-r}}\) 后,
我们可类似地, 从后向前地确定
\(x_{k_r}\), \(x_{k_{r-1}}\), \(\dots\), \(x_{k_1}\):
\begin{align*}
    x_{k_i} & =
    \frac{1}{a'_{i,k_i}}
    \left(
    b'_i
    - \sum_{1 \leq v \leq n-r} {a'_{i,\ell_v} x_{\ell_v}}
    - \sum_{i < w \leq r} {a'_{i,k_w} x_{k_w}}
    \right),
    \quad
    \text{\(i = r\), \(r - 1\), \(\dots\), \(1\)}.
\end{align*}
(这里, 注意, 若 \(i = r\), 则
\begin{align*}
    - \sum_{i < w \leq r} {a'_{i,k_w} x_{k_w}}
\end{align*}
不出现, 或被认为是 \(0\).)
注意, \(x_{\ell_1}\), \(\dots\), \(x_{\ell_{n-r}}\)
可取任何数.
则方程组~\eqref{eq:sole-5} 有不唯一的解.
则方程组~\eqref{eq:sole-4} 也有不唯一的解.

注意, \(r > n\) 是不可能的,
因为 \(r \leq k_r \leq n\).
(注意, \(1 \leq k_1 < k_2 < \dots < k_r \leq n\),
故 \(n \geq k_r \geq k_{r-1} + 1 \geq \dots \geq r_1 + r-1
\geq 1 + r - 1 = r\).)

综上, 我们有

\begin{theorem}
    考虑线性方程组~\eqref{eq:sole-4}.
    通过若干次 ``加一个方程的倍于另一个方程'' 的行为,
    我们可变方程组~\eqref{eq:sole-4}
    为方程组~\eqref{eq:sole-5},
    且它们有相同的解.

    (1)
    若 \(b'_{r+1} \neq 0\),
    则方程组~\eqref{eq:sole-4} 无解.

    (2)
    若 \(b'_{r+1} = 0\)
    (或 ``\(0 = b'_{r+1}\)'' 不出现),
    则方程组~\eqref{eq:sole-4} 有解.
    进一步地,
    若 \(r = n\),
    则方程组~\eqref{eq:sole-4} 有唯一的解;
    若 \(r < n\),
    则方程组~\eqref{eq:sole-4} 有不唯一的解.
    (\(r > n\) 是不可能的.)
\end{theorem}

我们考虑 REF 阵的秩.

\begin{theorem}
    设 \(F\) 是 REF 阵.
    设 \(F\) 有 \(r\) 个非零行.
    则 \(\operatorname{rank} {(F)} = r\).
\end{theorem}

\begin{proof}
    我们知道,
    \(\operatorname{pivot} {(F; 1)}
    < \operatorname{pivot} {(F; 2)}
    < \dots
    < \operatorname{pivot} {(F; r)}
    < \infty\),
    且对高于 \(r\) 的 \(i\),
    \(F\) 的行~\(i\) 是 \(0\).
    记 \(k_i = \operatorname{pivot} {(F; i)}\)
    (\(i = 1\), \(2\), \(\dots\), \(r\)).
    则 \([F]_{i,k_i} \neq 0\).
    则 \(F\) 有行列式非零的 \(r\)~级子阵, 因为
    \begin{align*}
        \det {\left(
            F\binom{1,2,\dots,r}{k_1,k_2,\dots,k_r}
            \right)}
        = {}    &
        \det {
            \begin{bmatrix}
                [F]_{1,k_1} & [F]_{1,k_2} & \cdots & [F]_{1,k_r} \\
                0           & [F]_{2,k_2} & \cdots & [F]_{2,k_r} \\
                \vdots      & \vdots      & \ddots & \vdots      \\
                0           & 0           & \cdots & [F]_{r,k_r} \\
            \end{bmatrix}
        }
        \\
        = {}    &
        [F]_{1,k_1} \dots [F]_{r,k_r}
        \\
        \neq {} &
        0.
    \end{align*}
    然后, \(F\) 没有行列式非零的 \(r + 1\)~级子阵.
    若 \(F\) 没有 \(r + 1\)~级子阵,
    则 \(F\) 当然没有行列式非零的 \(r + 1\)~级子阵.
    设 \(F\) 有 \(r + 1\)~级子阵.
    那么, 在 \(F\) 的任何 (不同的) \(r + 1\) 行,
    必存在一行, 其是 \(0\).
    (因为当 \(u\) 不等于 \(1\), \(2\), \(\dots\), \(r\)
    的任何一个时,
    \(F\) 的行~\(u\) 是 \(0\).)
    则 \(F\) 的任何 \(r + 1\)~级子阵必有一行, 其是 \(0\).
    则它的行列式为 \(0\).
    则 \(\operatorname{rank} {(F)} = r\).
\end{proof}

于是, 若我们用行的倍加变 \(A\) 为 REF 阵 \(F\),
则 \(A\)~的秩是 \(F\) 的非零行的数目,
因为倍加不改变秩.

\begin{theorem}
    设 \(A\) 是 \(m \times n\)~阵.
    设 \(B\) 是 \(m \times 1\)~阵.
    作 \(m \times (n + 1)\)~阵 \(G\) 如下:
    \begin{align*}
        [G]_{i,j}
        = \begin{cases}
              [A]_{i,j}, & j \leq n;  \\
              [B]_{i,1}, & j = n + 1.
          \end{cases}
    \end{align*}

    (1)
    设用若干次行的倍加,
    我们可变 \(G\) 为 REF 阵
    \begin{align*}
        G' =
        \begin{bmatrix}
            {}         & [G']_{1,k_1} & \cdots & [G']_{1,k_2} & \cdots     & [G']_{1,k_r} & \cdots & [G']_{1,n+1}   \\
            {}         & 0            & \cdots & [G']_{2,k_2} & \cdots     & [G']_{2,k_r} & \cdots & [G']_{2,n+1}   \\
            \mathsf{O} & \vdots       & {}     & \vdots       & {}         & \vdots       & \cdots & \vdots         \\
            {}         & 0            & \cdots & 0            & \cdots     & [G']_{r,k_r} & \cdots & [G']_{r,n+1}   \\
            {}         & 0            & \cdots & 0            & \cdots     & 0            & \cdots & [G']_{r+1,n+1} \\
            {}         & {}           & {}     & {}           & \mathsf{O} & {}           & {}     & {}             \\
        \end{bmatrix},
    \end{align*}
    其中,
    \(\mathsf{O}\) 表示可能出现的 \(0\),
    \(1 \leq k_1 < k_2 < \dots < k_r \leq n\),
    且 \(\operatorname{pivot} {(G'; i)} = k_i\)
    (\(i = 1\), \(2\), \(\dots\), \(r\)).
    (\([G']_{r+1,n+1}\) 可能不出现.)
    则由 \(G'\) 的前 \(n\)~列作成的
    \(m \times n\)~阵 \(A'\)
    也是 REF 阵,
    且这些行的倍加可变 \(A\) 为 \(A'\).

    (2)
    \(\operatorname{rank} {(A)}
    = \operatorname{rank} {(A')} = r\),
    且
    \(\operatorname{rank} {(G)} = \operatorname{rank} {(G')}\).

    (3)
    \([G']_{r+1,n+1} \neq 0\) 对应
    \(\operatorname{rank} {(G')} > \operatorname{rank} {(A')}\),
    即
    \(\operatorname{rank} {(G)} > \operatorname{rank} {(A)}\).

    (4)
    \([G']_{r+1,n+1} = 0\)
    (或 ``\([G']_{r+1,n+1}\)'' 不出现) 对应
    \(\operatorname{rank} {(G')} = \operatorname{rank} {(A')}\),
    即
    \(\operatorname{rank} {(G)} = \operatorname{rank} {(A)}\).
\end{theorem}

\begin{proof}
    % (1)
    % 我们知道, 我们可用若干次行的倍加,
    % 变 \(G\) 为 REF 阵 \(G'\).
    % 则存在非负整数 \(u\), 使
    % \begin{align*}
    %     \operatorname{pivot} {(G'; 1)}
    %     < \operatorname{pivot} {(G'; 2)}
    %     < \dots
    %     < \operatorname{pivot} {(G'; u-1)}
    %     < \operatorname{pivot} {(G'; u)}
    %     < \infty,
    % \end{align*}
    % 且对 \(i > u\),
    % \(G'\) 的行~\(i\) 是 \(0\).
    % 记 \(\operatorname{pivot} {(G'; i)} = \ell_i\)
    % (\(i = 1\), \(2\), \(\dots\), \(u\)).

    % 若 \(\ell_u = n + 1\),
    % 则 \(G'\) 形如
    % \begin{align*}
    %     \begin{bmatrix}
    %         {}         & [G']_{1,\ell_1} & \cdots & [G']_{1,\ell_2} & \cdots     & [G']_{1,\ell_{u-1}}   & \cdots & [G']_{1,n+1}   \\
    %         {}         & 0               & \cdots & [G']_{2,\ell_2} & \cdots     & [G']_{2,\ell_{u-1}}   & \cdots & [G']_{2,n+1}   \\
    %         \mathsf{O} & \vdots          & {}     & \vdots          & {}         & \vdots                & \cdots & \vdots         \\
    %         {}         & 0               & \cdots & 0               & \cdots     & [G']_{u-1,\ell_{u-1}} & \cdots & [G']_{u-1,n+1} \\
    %         {}         & 0               & \cdots & 0               & \cdots     & 0                     & \cdots & [G']_{u,n+1}   \\
    %         {}         & {}              & {}     & {}              & \mathsf{O} & {}                    & {}     & {}             \\
    %     \end{bmatrix},
    % \end{align*}
    % 其中, \(\mathsf{O}\) 表示可能出现的 \(0\), 下同.
    % 记 \(r = u - 1\),
    % 且 \(k_i = \ell_i\)
    % (\(i = 1\), \(2\), \(\dots\), \(u - 1\)).

    % 若 \(\ell_u \leq n\),
    % 则 \(G'\) 形如
    % \begin{align*}
    %     \begin{bmatrix}
    %         {}         & [G']_{1,\ell_1} & \cdots & [G']_{1,\ell_2} & \cdots     & [G']_{1,\ell_u} & \cdots & [G']_{1,n+1} \\
    %         {}         & 0               & \cdots & [G']_{2,\ell_2} & \cdots     & [G']_{2,\ell_u} & \cdots & [G']_{2,n+1} \\
    %         \mathsf{O} & \vdots          & {}     & \vdots          & {}         & \vdots          & \cdots & \vdots       \\
    %         {}         & 0               & \cdots & 0               & \cdots     & [G']_{u,\ell_u} & \cdots & [G']_{u,n+1} \\
    %         {}         & {}              & {}     & {}              & \mathsf{O} & {}              & {}     & {}           \\
    %     \end{bmatrix}
    % \end{align*}
    % 记 \(r = u\),
    % 且 \(k_i = \ell_i\)
    % (\(i = 1\), \(2\), \(\dots\), \(u\)).

    % 于是, 统一地, \(G'\) 形如
    % \begin{align*}
    %     \begin{bmatrix}
    %         {}         & [G']_{1,k_1} & \cdots & [G']_{1,k_2} & \cdots     & [G']_{1,k_r} & \cdots & [G']_{1,n+1}   \\
    %         {}         & 0            & \cdots & [G']_{2,k_2} & \cdots     & [G']_{2,k_r} & \cdots & [G']_{2,n+1}   \\
    %         \mathsf{O} & \vdots       & {}     & \vdots       & {}         & \vdots       & \cdots & \vdots         \\
    %         {}         & 0            & \cdots & 0            & \cdots     & [G']_{r,k_r} & \cdots & [G']_{r,n+1}   \\
    %         {}         & 0            & \cdots & 0            & \cdots     & 0            & \cdots & [G']_{r+1,n+1} \\
    %         {}         & {}           & {}     & {}           & \mathsf{O} & {}           & {}     & {}             \\
    %     \end{bmatrix},
    % \end{align*}
    % 其中,
    % \(1 \leq k_1 < k_2 < \dots < k_r \leq n\),
    % 且 \([G']_{i,k_i} \neq 0\)
    % (\(i = 1\), \(2\), \(\dots\), \(r\)).
    % (\([G']_{r+1,n+1}\) 可能不存在.)

    (1)
    设 \(A'\) 是由 \(G'\) 的前 \(n\) 列作成的 \(m \times n\)~阵.
    则
    \(\operatorname{pivot} {(A'; i)} = \operatorname{pivot} {(G'; i)}\)
    (\(i = 1\), \(2\), \(\dots\), \(r\)),
    且对 \(i > r\),
    \(A'\) 的行~\(i\) 是 \(0\).
    则 \(A'\) 是 REF 阵.

    设形如 \(E(m; p, q; s)\) (\(p \neq q\)) 的阵
    \(E_1\), \(E_2\), \(\dots\), \(E_w\) 使
    \begin{align*}
        G' = E_w (\dots (E_2 (E_1 G))).
    \end{align*}
    记 \(P = E_w \dots E_2 E_1\).
    则 \(G' = P G\).
    设 \(G\) 的列
    \(1\), \(2\), \(\dots\), \(n\), \(n+1\)
    是 \(g_1\), \(g_2\), \(\dots\), \(g_n\), \(g_{n+1}\).
    注意, \(g_j\) 也是 \(A\) 的列~\(j\)
    (\(j = 1\), \(2\), \(\dots\), \(n\)).
    则
    \begin{align*}
        G' = [P g_1, P g_2, \dots, P g_n, P g_{n+1}].
    \end{align*}
    则
    \begin{align*}
        A'
        = [P g_1, P g_2, \dots, P g_n]
        = P [g_1, g_2, \dots, g_n]
        = P A.
    \end{align*}
    则
    \begin{align*}
        A' = E_w (\dots (E_2 (E_1 A))).
    \end{align*}

    (2)
    倍加不改变秩.
    注意, \(A'\) 有 \(r\)~个非零行.

    (3) (4)
    若 \([G']_{r+1,n+1} \neq 0\),
    则 \(G'\) 有 \(r + 1\)~个非零行.
    则 \(\operatorname{rank} {(G')} > \operatorname{rank} {(A')}\).

    若 \([G']_{r+1,n+1} = 0\)
    (或 ``\([G']_{r+1,n+1}\)'' 不出现),
    则 \(G'\) 有 \(r\)~个非零行.
    则 \(\operatorname{rank} {(G')} = \operatorname{rank} {(A')}\).

    若
    \(\operatorname{rank} {(G')} > \operatorname{rank} {(A')}\),
    则 \([G']_{r+1,n+1}\) 不能是 \(0\)
    (若不然, 则
    \(\operatorname{rank} {(G')} = \operatorname{rank} {(A')}\);
    这是矛盾).

    若
    \(\operatorname{rank} {(G')} = \operatorname{rank} {(A')}\),
    则 ``\([G']_{r+1,n+1}\)'' 不出现,
    或 \([G']_{r+1,n+1} = 0\)
    (若不然, 则
    \(\operatorname{rank} {(G')} > \operatorname{rank} {(A')}\);
    这是矛盾).
\end{proof}

结合前面的讨论, 我们注意,
我们其实再证明了如下定理:

\begin{theorem}
    设 \(A\) 是 \(m \times n\)~阵.
    设 \(B\) 是 \(m \times 1\)~阵.
    作 \(m \times (n+1)\)~阵 \(G\),
    其中,
    \begin{align*}
        [G]_{i,j}
        = \begin{cases}
              [A]_{i,j}, & j \leq n;  \\
              [B]_{i,1}, & j = n + 1.
          \end{cases}
    \end{align*}

    (1)
    若 \(\operatorname{rank} {(G)}
    = \operatorname{rank} {(A)}\),
    则线性方程组 \(A X = B\) 有解.
    进一步地, 若 \(\operatorname{rank} {(A)} < n\),
    则 \(AX = B\) 的解不唯一;
    若 \(\operatorname{rank} {(A)} = n\),
    则 \(AX = B\) 的解唯一.

    (2)
    若 \(\operatorname{rank} {(G)}
    > \operatorname{rank} {(A)}\),
    则线性方程组 \(A X = B\) 无解.

    注意, \(\operatorname{rank} {(G)}\)
    不可能小于 \(\operatorname{rank} {(A)}\).
\end{theorem}

于是, 我们能用 REF 阵判断线性方程组是否有解,
判断线性方程组的解是否唯一 (若解存在),
解线性方程组 (若解存在),
且求阵的秩.
用行的倍加变阵为 REF 阵或许比%
求它的子阵的行列式简单.

\begin{example}
    解线性方程组
    \begin{equation}
        \left\{
        \begin{aligned}
            0 x_1 + 1 x_2 + 1 x_3 + 1 x_4 + 1 x_5 & = 17, \\
            1 x_1 + 0 x_2 + 1 x_3 + 1 x_4 + 1 x_5 & = 15, \\
            1 x_1 + 1 x_2 + 0 x_3 + 1 x_4 + 1 x_5 & = 14, \\
            1 x_1 + 1 x_2 + 1 x_3 + 0 x_4 + 1 x_5 & = 12, \\
            1 x_1 + 1 x_2 + 1 x_3 + 1 x_4 + 0 x_5 & = 10.
        \end{aligned}
        \right.
        \label{eq:sole-6}
    \end{equation}

    记
    \begin{align*}
        A =
        \begin{bmatrix}
            0 & 1 & 1 & 1 & 1 \\
            1 & 0 & 1 & 1 & 1 \\
            1 & 1 & 0 & 1 & 1 \\
            1 & 1 & 1 & 0 & 1 \\
            1 & 1 & 1 & 1 & 0 \\
        \end{bmatrix},
        \quad
        X =
        \begin{bmatrix}
            x_1 \\
            x_2 \\
            x_3 \\
            x_4 \\
            x_5 \\
        \end{bmatrix},
        \quad
        B =
        \begin{bmatrix}
            17 \\
            15 \\
            14 \\
            12 \\
            10 \\
        \end{bmatrix}.
    \end{align*}
    则方程组~\eqref{eq:sole-6} 相当于 \(A X = B\).

    线性方程组~\eqref{eq:sole-6}
    对应 \(5 \times (5 + 1)\)~阵
    \begin{align*}
        G = [A, B] =
        \begin{bmatrix}
            0 & 1 & 1 & 1 & 1 & 17 \\
            1 & 0 & 1 & 1 & 1 & 15 \\
            1 & 1 & 0 & 1 & 1 & 14 \\
            1 & 1 & 1 & 0 & 1 & 12 \\
            1 & 1 & 1 & 1 & 0 & 10 \\
        \end{bmatrix}.
    \end{align*}
    由前面的例, 我们可用行的倍加, 变 \(G\) 为
    \begin{align*}
        G' =
        \begin{bmatrix}
            4 & 0  & 0  & 0  & 0  & 0  \\
            0 & -1 & 0  & 0  & 0  & -2 \\
            0 & 0  & -1 & 0  & 0  & -3 \\
            0 & 0  & 0  & -1 & 0  & -5 \\
            0 & 0  & 0  & 0  & -1 & -7 \\
        \end{bmatrix}.
    \end{align*}
    设 \(A'\) 是由 \(G'\) 的前 \(5\) 列作成的阵.
    则
    \(\operatorname{rank} {(G)}
    = \operatorname{rank} {(G')}
    = \operatorname{rank} {(A')}
    = \operatorname{rank} {(A)}
    = 5\).
    故方程组~\eqref{eq:sole-6} 有解.
    因为 \(A\)~的秩等于 \(A\)~的列数,
    方程组~\eqref{eq:sole-6} 有唯一的解.

    \(G'\) 对应线性方程组
    \begin{equation*}
        \left\{
        \begin{aligned}
            4 x_1 & = 0,  \\
            - x_2 & = -2, \\
            - x_3 & = -3, \\
            - x_4 & = -5, \\
            - x_5 & = -7.
        \end{aligned}
        \right.
    \end{equation*}
    它与方程组~\eqref{eq:sole-6} 有相同的解.
    由此, 我们可解得
    \(x_1 = 0\), \(x_2 = 2\), \(x_3 = 3\),
    \(x_4 = 5\), \(x_5 = 7\).

    由上面的计算,
    我们看出,
    \(\det {(A)} = \det {(A')}
    = 4 \cdot (-1)^4 = 4\).
    于是, 理论地, 我们也可用 Cramer 公式%
    解线性方程组~\eqref{eq:sole-6}.
    不过, 计算行列式不是简单的事情:
    若我们用 Cramer 公式~1,
    \begin{align*}
        X = (\det {(A)})^{-1} \operatorname{adj} {(A)}\, B,
    \end{align*}
    则我们要计算 \(5 \cdot 5 = 25\)~个 \(4\)~级阵的行列式;
    若我们用 Cramer 公式~2,
    \begin{align*}
        x_i = \frac{\det {(A \{i, B\})}}{\det {(A)}},
        \quad
        i = 1, 2, 3, 4, 5
    \end{align*}
    (其中, \(A \{i, B\}\) 是以
    \(B\) 代阵~\(A\) 的列~\(i\) 后得到的阵),
    则我们要计算 \(5\)~个 \(5\)~级阵的行列式.
\end{example}

\begin{example}
    设 \(k\) 是参数.
    考虑由 \(3\)~个%
    关于 \(x\), \(y\) 的
    \(2\)~元 \(1\)~次方程作成的方程组
    \begin{equation}
        \left\{
        \begin{alignedat}{4}
            1 x              & {} + {} & k y & {} = {} & 1,             \\
            k x              & {} + {} & 1 y & {} = {} & 1,             \\
            (-k^2 + k + 1) x & {} + {} & 1 y & {} = {} & k^2 - 2 k + 2, \\
        \end{alignedat}
        \right.
        \label{eq:sole-7}
    \end{equation}
    若方程组~\eqref{eq:sole-7} 有解,
    \(k\) 应适合什么条件?
    反过来, 当 \(k\) 适合什么条件时,
    方程组~\eqref{eq:sole-7} 有解?
    当方程组~\eqref{eq:sole-7} 有解时,
    求它的解.

    记
    \begin{align*}
        A =
        \begin{bmatrix}
            1        & k \\
            k        & 1 \\
            -k^2+k+1 & 1 \\
        \end{bmatrix},
        \quad
        X =
        \begin{bmatrix}
            x \\
            y \\
        \end{bmatrix},
        \quad
        B =
        \begin{bmatrix}
            1        \\
            1        \\
            k^2-2k+2 \\
        \end{bmatrix}.
    \end{align*}
    则方程组~\eqref{eq:sole-7} 相当于 \(A X = B\).

    线性方程组~\eqref{eq:sole-7}
    对应 \(3 \times (2 + 1)\)~阵
    \begin{align*}
        G = [A, B] =
        \begin{bmatrix}
            1        & k & 1        \\
            k        & 1 & 1        \\
            -k^2+k+1 & 1 & k^2-2k+2 \\
        \end{bmatrix}.
    \end{align*}
    加行~\(1\) 的 \(-k\)~倍于行~\(2\),
    且加行~\(1\) 的 \(k^2 - k - 1\)~倍于行~\(3\),
    得
    \begin{align*}
        \begin{bmatrix}
            1 & k            & 1           \\
            0 & 1-k^2        & 1-k         \\
            0 & (1-k^2)(1-k) & (1-k)(1-2k) \\
        \end{bmatrix}.
    \end{align*}
    加行~\(2\) 的 \(k-1\)~倍于行~\(3\),
    得
    \begin{align*}
        \begin{bmatrix}
            1 & k     & 1     \\
            0 & 1-k^2 & 1-k   \\
            0 & 0     & k^2-k \\
        \end{bmatrix}
        = H.
    \end{align*}
    注意, \(H\) 不一定是 REF 阵,
    因为 \(k\) 是参数.
    我们要分类讨论.

    设 \(1 - k^2 \neq 0\).
    则 \(H\) 是 REF 阵.
    设 \(G' = H\).
    设 \(A'\) 是由 \(G'\) 的前 \(2\) 列作成的阵.
    则 \(\operatorname{rank} {(A)}
    = \operatorname{rank} {(A')} = 2\).
    注意,
    \begin{align*}
        \operatorname{rank} {(G)}
        = {} &
        \operatorname{rank} {(G')}
        \\
        = {} &
        \begin{cases}
            2, & \text{\(k^2-k = 0\), 且 \(1-k^2 \neq 0\)};   \\
            3, & \text{\(k^2-k \neq 0\), 且 \(1-k^2 \neq 0\)}
        \end{cases}
        \\
        = {} &
        \begin{cases}
            2, & k = 0;                                              \\
            3, & \text{\(k \neq 0\) 且 \(k \neq 1\) 且 \(k \neq -1\)}.
        \end{cases}
    \end{align*}
    所以, 当 \(1 - k^2 \neq 0\) 时,
    若方程组~\eqref{eq:sole-7} 有解, 必
    \(\operatorname{rank} {(G)} = \operatorname{rank} {(A)}\),
    即 \(k = 0\).
    反过来, 当 \(k = 0\) 时,
    \(\operatorname{rank} {(G)} = \operatorname{rank} {(A)}\),
    故方程组~\eqref{eq:sole-7} 有解.
    因为 \(A\)~的秩等于 \(A\)~的列数,
    则当 \(k = 0\) 时,
    方程组~\eqref{eq:sole-7} 有唯一的解.

    进一步地, 当 \(k = 0\) 时,
    \(G'\) 对应线性方程组
    \begin{equation*}
        \left\{
        \begin{aligned}
            x & = 1, \\
            y & = 1, \\
            0 & = 0.
        \end{aligned}
        \right.
    \end{equation*}
    它与方程组~\eqref{eq:sole-7}
    (其中, \(k = 0\)) 有相同的解.
    则 \(x = 1\), \(y = 1\).

    设 \(1 - k^2 = 0\).
    则 \(k = 1\), 或 \(k = -1\).
    则
    \begin{align*}
        H =
        \begin{bmatrix}
            1 & k & 1   \\
            0 & 0 & 1-k \\
            0 & 0 & 1-k \\
        \end{bmatrix}.
    \end{align*}
    加行~\(2\) 的 \(-1\)~倍于行~\(3\),
    得
    \begin{align*}
        \begin{bmatrix}
            1 & k & 1   \\
            0 & 0 & 1-k \\
            0 & 0 & 0   \\
        \end{bmatrix}
        = G'.
    \end{align*}
    \(G'\) 是 REF 阵.
    设 \(A'\) 是由 \(G'\) 的前 \(2\) 列作成的阵.
    则 \(\operatorname{rank} {(A)}
    = \operatorname{rank} {(A')} = 1\).
    注意,
    \begin{align*}
        \operatorname{rank} {(G)}
        = \operatorname{rank} {(G')}
        = \begin{cases}
              1, & k = 1;  \\
              2, & k = -1.
          \end{cases}
    \end{align*}
    所以, 当 \(1 - k^2 = 0\) 时,
    若方程组~\eqref{eq:sole-7} 有解,
    必
    \(\operatorname{rank} {(G)} = \operatorname{rank} {(A)}\),
    即 \(k = 1\).
    反过来, 当 \(k = 1\) 时,
    \(\operatorname{rank} {(G)} = \operatorname{rank} {(A)}\),
    故方程组~\eqref{eq:sole-7} 有解.
    因为 \(A\)~的秩小于 \(A\)~的列数,
    则当 \(k = 1\) 时,
    方程组~\eqref{eq:sole-7} 有不唯一的解.

    进一步地, 当 \(k = 1\) 时, \(G'\) 对应线性方程组
    \begin{equation*}
        \left\{
        \begin{aligned}
            x + y & = 1, \\
            0     & = 0, \\
            0     & = 0.
        \end{aligned}
        \right.
    \end{equation*}
    它与方程组~\eqref{eq:sole-7}
    (其中, \(k = 1\)) 有相同的解.
    则 \(x = 1 - c\), \(y = c\),
    其中, \(c\) 是任何数.
\end{example}

前面, 我们用 REF 阵得到了线性方程组的一些理论.
我们说, 我们还能得到更多.

\begin{definition}
    设 \(A\) 是 \(m \times n\)~阵.
    设存在不高于 \(m\) 的非负整数 \(r\), 使:

    (1)
    对任何不高于 \(r\) 的正整数 \(i\),
    \begin{align*}
        \operatorname{pivot} {(A; i)} < \infty;
    \end{align*}

    (2)
    对任何低于 \(r\) 的正整数 \(i\),
    \begin{align*}
        \operatorname{pivot} {(A; i)}
        < \operatorname{pivot} {(A; i + 1)};
    \end{align*}

    (3)
    对任何高于 \(r\) 且不高于 \(m\) 的整数 \(i\),
    \begin{align*}
        \operatorname{pivot} {(A; i)} = \infty;
    \end{align*}

    (4)
    对任何不高于 \(r\) 的正整数 \(i\),
    与任何不等于 \(i\),
    且不高于 \(m\) 的正整数 \(\ell\),
    \begin{align*}
        [A]_{\ell,\operatorname{pivot} {(A; i)}} = 0.
    \end{align*}

    我们说, \(A\) 是 \emph{SREF}
    (\angla{special row-echelon form}) \emph{阵}.
\end{definition}

注意, SREF 阵是 REF 阵,
但 REF 阵不一定是 SREF 阵.

\begin{example}
    设
    \begin{align*}
        G =
        \begin{bmatrix}
            2 & 4 & 94 \\
            1 & 1 & 35 \\
        \end{bmatrix}.
    \end{align*}
    注意, \(G\) 不是 REF 阵, 且 \(G\) 不是 SREF 阵.

    对 \(G\),
    我们加行~\(1\) 的 \(-1/2\) 倍于行~\(2\),
    得
    \begin{align*}
        G_1 =
        \begin{bmatrix}
            2 & 4  & 94  \\
            0 & -1 & -12 \\
        \end{bmatrix}.
    \end{align*}
    注意, \(G_1\) 是 REF 阵, 但 \(G_1\) 不是 SREF 阵,
    因为 \(1 \neq \operatorname{pivot} {(G_1; 2)} = 2\),
    但 \([G_1]_{1,2} = 4 \neq 0\).

    对 \(G_1\),
    我们加行~\(2\) 的 \(4\) 倍于行~\(1\),
    得
    \begin{align*}
        G_2 =
        \begin{bmatrix}
            2 & 0  & 46  \\
            0 & -1 & -12 \\
        \end{bmatrix}.
    \end{align*}
    注意, \(G_2\) 是 REF 阵, 且 \(G_2\) 是 SREF 阵.
\end{example}

我们知道, 我们可用行的倍加,
变一个阵为 REF 阵.
进一步地, 我们有

\begin{theorem}
    设 \(A\) 是 \(m \times n\)~阵.
    则我们可用行的倍加,
    变 \(A\) 为某 \(m \times n\) SREF 阵 \(D\).
\end{theorem}

\begin{proof}
    我们可用行的倍加,
    变 \(A\) 为某 \(m \times n\) REF 阵 \(B\).
    则存在非负整数 \(r\), 使
    \begin{align*}
        \operatorname{pivot} {(B; 1)}
        < \operatorname{pivot} {(B; 2)}
        < \dots
        < \operatorname{pivot} {(B; r-1)}
        < \operatorname{pivot} {(B; r)}
        < \infty,
    \end{align*}
    且对 \(i > r\),
    \(B\) 的行~\(i\) 是 \(0\).
    记 \(\operatorname{pivot} {(B; i)} = k_i\)
    (\(i = 1\), \(2\), \(\dots\), \(r\)).

    注意, 若 \(\ell > i\), 则 \([B]_{\ell,k_i} = 0\).
    于是, 若 \(i \leq 1\),
    且 \(\ell \neq i\),
    则 \([B]_{\ell,k_i} = 0\).

    对 \(B\),
    我们加行~\(2\) 的
    \(-[B]_{1,k_2} / [B]_{2,k_2}\)~倍于行~\(1\),
    得阵 \(B_2\).
    于是, 若 \(\ell \neq 2\), 则 \([B_2]_{\ell,k_2} = 0\).
    注意, \(\operatorname{pivot} {(B_2; i)}
    = \operatorname{pivot} {(B; i)}\),
    且若 \(j < k_1\),
    则 \([B_2]_{i,j} = [B]_{i,j}\).
    于是, 若 \(i \leq 2\),
    且 \(\ell \neq i\),
    则 \([B_2]_{\ell,k_i} = 0\).

    对 \(B_2\),
    我们加行~\(3\) 的
    \(-[B_2]_{1,k_3} / [B_2]_{3,k_3}\)~倍于行~\(1\),
    且加行~\(3\) 的
    \(-[B_2]_{2,k_3} / [B_2]_{3,k_3}\)~倍于行~\(2\),
    得阵 \(B_3\).
    于是, 若 \(\ell \neq 3\), 则 \([B_3]_{\ell,k_3} = 0\).
    注意, \(\operatorname{pivot} {(B_3; i)}
    = \operatorname{pivot} {(B_2; i)}
    = \operatorname{pivot} {(B; i)}\),
    且若 \(j < k_2\),
    则 \([B_3]_{i,j} = [B_2]_{i,j}\).
    于是, 若 \(i \leq 3\),
    且 \(\ell \neq i\),
    则 \([B_3]_{\ell,k_i} = 0\).

    ……
    最后, 用若干次行的倍加,
    我们得 \(m \times n\)~阵 \(D\),
    使 \(\operatorname{pivot} {(D; i)}
    = \operatorname{pivot} {(B; i)}\).
    并且, 若 \(i \leq r\),
    且 \(\ell \neq i\),
    则 \([D]_{\ell,k_i} = 0\).
    则 \(D\) 是 SREF 阵.
\end{proof}

\begin{theorem}
    设 \(A\) 是 \(m \times n\)~阵.
    设 \(B\) 是 \(m \times 1\)~阵.
    作 \(m \times (n + 1)\)~阵
    \(G = [A, B]\).
    则我们可用行的倍加,
    变 \(G\) 为 REF 阵 \(G'\),
    且由 \(G'\) 的前 \(n\)~列作成的
    \(m \times n\)~阵 \(A'\)
    是 SREF 阵.
\end{theorem}

\begin{proof}
    我们知道, 我们可用若干次行的倍加,
    变 \(G\) 为 SREF 阵 \(G'\).
    则存在非负整数 \(u\), 使
    \begin{align*}
        \operatorname{pivot} {(G'; 1)}
        < \operatorname{pivot} {(G'; 2)}
        < \dots
        < \operatorname{pivot} {(G'; u-1)}
        < \operatorname{pivot} {(G'; u)}
        < \infty,
    \end{align*}
    且对 \(i > u\),
    \(G'\) 的行~\(i\) 是 \(0\).
    记 \(\operatorname{pivot} {(G'; i)} = p_i\)
    (\(i = 1\), \(2\), \(\dots\), \(u\)).
    则当 \(i \leq u\),
    且 \(\ell \neq i\) 时,
    \([G']_{\ell,k_i} = 0\).

    若 \(p_u = n + 1\),
    则 \(G'\) 形如
    \begin{align*}
        \begin{bmatrix}
            {}         & [G']_{1,p_1} & \cdots & 0            & \cdots     & 0                  & \cdots & 0            \\
            {}         & 0            & \cdots & [G']_{2,p_2} & \cdots     & 0                  & \cdots & 0            \\
            \mathsf{O} & \vdots       & {}     & \vdots       & {}         & \vdots             & \cdots & \vdots       \\
            {}         & 0            & \cdots & 0            & \cdots     & [G']_{u-1,p_{u-1}} & \cdots & 0            \\
            {}         & 0            & \cdots & 0            & \cdots     & 0                  & \cdots & [G']_{u,n+1} \\
            {}         & {}           & {}     & {}           & \mathsf{O} & {}                 & {}     & {}           \\
        \end{bmatrix},
    \end{align*}
    其中, \(\mathsf{O}\) 表示可能出现的 \(0\), 下同.
    记 \(r = u - 1\),
    且 \(k_i = p_i\)
    (\(i = 1\), \(2\), \(\dots\), \(u - 1\)).

    若 \(p_u \leq n\),
    则 \(G'\) 形如
    \begin{align*}
        \begin{bmatrix}
            {}         & [G']_{1,p_1} & \cdots & 0            & \cdots     & 0            & \cdots & [G']_{1,n+1} \\
            {}         & 0            & \cdots & [G']_{2,p_2} & \cdots     & 0            & \cdots & [G']_{2,n+1} \\
            \mathsf{O} & \vdots       & {}     & \vdots       & {}         & \vdots       & \cdots & \vdots       \\
            {}         & 0            & \cdots & 0            & \cdots     & [G']_{u,p_u} & \cdots & [G']_{u,n+1} \\
            {}         & {}           & {}     & {}           & \mathsf{O} & {}           & {}     & {}           \\
        \end{bmatrix}
    \end{align*}
    记 \(r = u\),
    且 \(k_i = p_i\)
    (\(i = 1\), \(2\), \(\dots\), \(u\)).

    于是, 统一地, \(G'\) 形如
    \begin{align*}
        \begin{bmatrix}
            {}         & [G']_{1,k_1} & \cdots & 0            & \cdots     & 0            & \cdots & [G']_{1,n+1}   \\
            {}         & 0            & \cdots & [G']_{2,k_2} & \cdots     & 0            & \cdots & [G']_{2,n+1}   \\
            \mathsf{O} & \vdots       & {}     & \vdots       & {}         & \vdots       & \cdots & \vdots         \\
            {}         & 0            & \cdots & 0            & \cdots     & [G']_{r,k_r} & \cdots & [G']_{r,n+1}   \\
            {}         & 0            & \cdots & 0            & \cdots     & 0            & \cdots & [G']_{r+1,n+1} \\
            {}         & {}           & {}     & {}           & \mathsf{O} & {}           & {}     & {}             \\
        \end{bmatrix},
    \end{align*}
    其中,
    \(1 \leq k_1 < k_2 < \dots < k_r \leq n\),
    且 \([G']_{i,k_i} \neq 0\)
    (\(i = 1\), \(2\), \(\dots\), \(r\)),
    且当 \(i > r + 1\) 时, \([G']_{i,n+1} = 0\),
    且当 \(i > r\), 且 \(j \leq n\) 时,
    \([G']_{i,j} = 0\),
    且当 \(i \leq r\), 且 \(\ell \neq i\) 时,
    \([G']_{\ell,k_i} = 0\).
    (注意, 若 \(r = m\), 则 \([G']_{r+1,n+1}\) 不出现.)

    设 \(A'\) 是由 \(G'\) 的前 \(n\)~列作成的
    \(m \times n\)~阵.
    则
    \(\operatorname{pivot} {(A'; i)} = \operatorname{pivot} {(G'; i)}\)
    (\(i = 1\), \(2\), \(\dots\), \(r\)),
    且对 \(i > r\),
    \(A'\) 的行~\(i\) 是 \(0\),
    且当 \(i \leq r\), 且 \(\ell \neq i\) 时,
    \([A']_{\ell,k_i} = [G']_{\ell,k_i} = 0\).
    则 \(A'\) 是 SREF 阵.
\end{proof}

再考虑线性方程组
\begin{equation}
    \left\{
    \begin{aligned}
        a_{1,1} x_1 + a_{1,2} x_2 + \dots + a_{1,n} x_n & = b_1, \\
        a_{2,1} x_1 + a_{2,2} x_2 + \dots + a_{2,n} x_n & = b_2, \\
                                                        & \dots, \\
        a_{m,1} x_1 + a_{m,2} x_2 + \dots + a_{m,n} x_n & = b_m.
    \end{aligned}
    \right.
    \label{eq:sole-8}
\end{equation}
它跟 \(m \times (n + 1)\)~阵
\begin{align*}
    G =
    \begin{bmatrix}
        a_{1,1} & a_{1,2} & \cdots & a_{1,n} & b_1    \\
        a_{2,1} & a_{2,2} & \cdots & a_{2,n} & b_2    \\
        \vdots  & \vdots  & {}     & \vdots  & \vdots \\
        a_{m,1} & a_{m,2} & \cdots & a_{m,n} & b_m    \\
    \end{bmatrix}
\end{align*}
对应.
我们知道, 我们可用行的倍加,
变 \(G\) 为 REF 阵
\begin{align*}
    G'' =
    \begin{bmatrix}
        {}         & [G'']_{1,k_1} & \cdots & 0             & \cdots     & 0             & \cdots & [G'']_{1,n+1}   \\
        {}         & 0             & \cdots & [G'']_{2,k_2} & \cdots     & 0             & \cdots & [G'']_{2,n+1}   \\
        \mathsf{O} & \vdots        & {}     & \vdots        & {}         & \vdots        & \cdots & \vdots          \\
        {}         & 0             & \cdots & 0             & \cdots     & [G'']_{r,k_r} & \cdots & [G'']_{r,n+1}   \\
        {}         & 0             & \cdots & 0             & \cdots     & 0             & \cdots & [G'']_{r+1,n+1} \\
        {}         & {}            & {}     & {}            & \mathsf{O} & {}            & {}     & {}              \\
    \end{bmatrix},
\end{align*}
其中,
\(1 \leq k_1 < k_2 < \dots < k_r \leq n\),
且 \([G'']_{i,k_i} \neq 0\)
(\(i = 1\), \(2\), \(\dots\), \(r\)),
且当 \(i > r + 1\) 时, \([G'']_{i,n+1} = 0\),
且当 \(i > r\), 且 \(j \leq n\) 时,
\([G'']_{i,j} = 0\),
且当 \(i \leq r\), 且 \(\ell \neq i\) 时,
\([G'']_{\ell,k_i} = 0\).
(注意, 若 \(r = m\), 则 \([G'']_{r+1,n+1}\) 不出现.)

为方便, 记 \([G'']_{i,n+1} = b''_i\).
则当 \(i > r + 1\) 时, \(b''_i = 0\).
记 \([G'']_{i,j} = a''_{i,j}\)
(若 \(j \leq n\)).
则当 \(i > r\), 且 \(j \leq n\) 时,
\(a''_{i,j} = 0\),
且当 \(i \leq r\), 且 \(\ell \neq i\) 时,
\(a''_{\ell,k_i} = 0\).
则 \(G''\) 形如
\begin{align*}
    \begin{bmatrix}
        {}         & a''_{1,k_1} & \cdots & 0           & \cdots     & 0           & \cdots & b''_1     \\
        {}         & 0           & \cdots & a''_{2,k_2} & \cdots     & 0           & \cdots & b''_2     \\
        \mathsf{O} & \vdots      & {}     & \vdots      & {}         & \vdots      & \cdots & \vdots    \\
        {}         & 0           & \cdots & 0           & \cdots     & a''_{r,k_r} & \cdots & b''_r     \\
        {}         & 0           & \cdots & 0           & \cdots     & 0           & \cdots & b''_{r+1} \\
        {}         & {}          & {}     & {}          & \mathsf{O} & {}          & {}     & {}        \\
    \end{bmatrix}.
\end{align*}
设从 \(1\), \(2\), \(\dots\), \(n\)
去除 \(k_1\), \(\dots\), \(k_r\) 后,
还剩 \(n - r\)~个数.
我们从小到大地叫这 \(n - r\)~个数为
\(\ell_1\), \(\dots\), \(\ell_{n-r}\).
(注意, 若 \(r = n\), 则
\(\ell_1\), \(\dots\), \(\ell_{n-r}\)
不出现.)
则 \(G''\) 对应线性方程组
\begin{equation}
    \left\{
    \begin{aligned}
        a''_{1,k_1} x_{k_1} + \sum_{1 \leq v \leq n-r} {a''_{1,\ell_v} x_{\ell_v}} & = b''_1,     \\
        a''_{2,k_2} x_{k_2} + \sum_{1 \leq v \leq n-r} {a''_{2,\ell_v} x_{\ell_v}} & = b''_2,     \\
                                                                                   & \dots,       \\
        a''_{r,k_r} x_{k_r} + \sum_{1 \leq v \leq n-r} {a''_{r,\ell_v} x_{\ell_v}} & = b''_r,     \\
        0                                                                          & = b''_{r+1}, \\
        0                                                                          & = 0,         \\
                                                                                   & \dots,       \\
        0                                                                          & = 0,
    \end{aligned}
    \right.
    \label{eq:sole-9}
\end{equation}
其中, \(1 \leq k_1 < k_2 < \dots < k_r \leq n\),
且 \(a''_{i,k_i} \neq 0\)
(\(i = 1\), \(2\), \(\dots\), \(r\)).
(注意, 若 \(r = n\), 则
\begin{align*}
    + \sum_{1 \leq v \leq n-r} {a''_{i,\ell_v} x_{\ell_v}},
    \quad
    i = 1, 2, \dots, r
\end{align*}
不出现, 或被认为是 \(0\).)
由前面的讨论, 方程组~\eqref{eq:sole-9}
与方程组~\eqref{eq:sole-8} 有相同的解.
(注意, 形式地, 方程组~\eqref{eq:sole-9}
比方程组~\eqref{eq:sole-5} 更简单.)

类似地, 设 \(b''_{r+1} \neq 0\).
则方程组~\eqref{eq:sole-9} 的方程
``\(0 = b''_{r+1}\)'' 是矛盾.
则方程组~\eqref{eq:sole-9} 无解.
则方程组~\eqref{eq:sole-8} 也无解.
(注意, 此时必 \(r < m\).)

类似地, 设 \(b''_{r+1} = 0\) (当 \(r < m\)),
或 ``\(0 = b''_{r+1}\)'' 不出现 (当 \(r = m\)).

若 \(r = n\),
则 \(k_i = i\) (\(i = 1\), \(2\), \(\dots\), \(n\)).
则方程组~\eqref{eq:sole-9} 形如
\begin{equation*}
    \left\{
    \begin{aligned}
        a''_{1,1} x_1 & = b''_1, \\
        a''_{2,2} x_2 & = b''_2, \\
                      & \dots,   \\
        a''_{n,n} x_n & = b''_n,
    \end{aligned}
    \right.
\end{equation*}
其中, \(a''_{i,i} \neq 0\).
这里, 我们未写形如 ``\(0 = 0\)'' 的方程 (若它们出现),
因为它们是恒等式,
且去除它们不影响方程组的解.
由此, 我们可直接地确定
\begin{align*}
    x_i = \frac{b''_i}{a''_{i,i}},
    \quad i = 1, 2, \dots, n.
\end{align*}
(注意, 形式地, 这儿的公式更简单.)
则方程组~\eqref{eq:sole-9} 有唯一的解.
则方程组~\eqref{eq:sole-8} 有唯一的解.

若 \(r < n\), 我们移含
\(x_{\ell_1}\), \(\dots\), \(x_{\ell_{n-r}}\) 的项到右侧,
得
\begin{equation*}
    \left\{
    \begin{aligned}
        a''_{1,k_1} x_{k_1} & = b''_1 - \sum_{1 \leq v \leq n-r} {a''_{1,\ell_v} x_{\ell_v}}, \\
        a''_{2,k_2} x_{k_2} & = b''_2 - \sum_{1 \leq v \leq n-r} {a''_{2,\ell_v} x_{\ell_v}}, \\
                            & \dots,                                                          \\
        a''_{r,k_r} x_{k_r} & = b''_r - \sum_{1 \leq v \leq n-r} {a''_{r,\ell_v} x_{\ell_v}},
    \end{aligned}
    \right.
\end{equation*}
其中, \(a''_{i,k_i} \neq 0\).
这里, 我们未写形如 ``\(0 = 0\)'' 的方程 (若它们出现),
因为它们是恒等式,
且去除它们不影响方程组的解.
选取 \(x_{\ell_1}\), \(\dots\), \(x_{\ell_{n-r}}\) 后,
我们可确定 \(x_{k_1}\), \(\dots\), \(x_{k_r}\):
\begin{align*}
    x_{k_i} =
    \frac{1}{a''_{i,k_i}}
    \left(
    b''_i
    - \sum_{1 \leq v \leq n-r} {a''_{i,\ell_v} x_{\ell_v}}
    \right),
    \quad
    i = 1, 2, \dots, r.
\end{align*}
(注意, 形式地, 这儿的公式更简单.)
注意, \(x_{\ell_1}\), \(\dots\), \(x_{\ell_{n-r}}\)
可取任何数.
则方程组~\eqref{eq:sole-9} 有不唯一的解.
则方程组~\eqref{eq:sole-8} 也有不唯一的解.

由此可见,
若我们变对应线性方程组~\eqref{eq:sole-8} 的
\(m \times (n + 1)\)~阵 \(G\) 为 REF 阵 \(G'\),
且进一步地要求,
\(G'\) 的前 \(n\)~列作成 \(m \times n\) SREF 阵,
则对应 \(G'\) 的线性方程组~\eqref{eq:sole-9} 的解%
可被直接地写出.
特别地, 我们可再得到

\begin{theorem}
    考虑线性方程组
    \begin{equation}
        \left\{
        \begin{aligned}
            a_{1,1} x_1 + a_{1,2} x_2 + \dots + a_{1,n} x_n & = b_1, \\
            a_{2,1} x_1 + a_{2,2} x_2 + \dots + a_{2,n} x_n & = b_2, \\
                                                            & \dots, \\
            a_{n,1} x_1 + a_{n,2} x_2 + \dots + a_{n,n} x_n & = b_n.
        \end{aligned}
        \right.
        \label{eq:sole-10}
    \end{equation}
    (注意, 方程数等于未知数数.)
    记
    \begin{align*}
        A =
        \begin{bmatrix}
            a_{1,1} & a_{1,2} & \cdots & a_{1,n} \\
            a_{2,1} & a_{2,2} & \cdots & a_{2,n} \\
            \vdots  & \vdots  & {}     & \vdots  \\
            a_{n,1} & a_{n,2} & \cdots & a_{n,n} \\
        \end{bmatrix},
        \quad
        B =
        \begin{bmatrix}
            b_1 \\ b_2 \\ \vdots \\ b_n
        \end{bmatrix}.
    \end{align*}

    (1)
    (Cramer 公式, 2)
    若 \(\det {(A)} \neq 0\),
    则方程组~\eqref{eq:sole-10} 有唯一的解
    \begin{align*}
        x_i = \frac{\det {(A \{i, B\})}}{\det {(A)}},
        \quad
        \text{\(i = 1\), \(2\), \(\dots\), \(n\)},
    \end{align*}
    其中, \(A \{i, B\}\) 是%
    以 \(B\) 代阵~\(A\) 的列~\(i\) 后得到的阵.

    (2)
    若 \(\det {(A)} = 0\),
    则方程组~\eqref{eq:sole-10} 要么无解,
    要么有不唯一的解.
\end{theorem}

\begin{proof}
    作 \(n \times (n + 1)\) 阵 \(G = [A, B]\).
    我们可用行的倍加, 变 \(G\) 为 REF 阵
    \begin{align*}
        G' =
        \begin{bmatrix}
            {}         & [G']_{1,k_1} & \cdots & 0            & \cdots     & 0            & \cdots & [G']_{1,n+1}   \\
            {}         & 0            & \cdots & [G']_{2,k_2} & \cdots     & 0            & \cdots & [G']_{2,n+1}   \\
            \mathsf{O} & \vdots       & {}     & \vdots       & {}         & \vdots       & \cdots & \vdots         \\
            {}         & 0            & \cdots & 0            & \cdots     & [G']_{r,k_r} & \cdots & [G']_{r,n+1}   \\
            {}         & 0            & \cdots & 0            & \cdots     & 0            & \cdots & [G']_{r+1,n+1} \\
            {}         & {}           & {}     & {}           & \mathsf{O} & {}           & {}     & {}             \\
        \end{bmatrix},
    \end{align*}
    其中,
    \(1 \leq k_1 < k_2 < \dots < k_r \leq n\),
    且 \([G']_{i,k_i} \neq 0\)
    (\(i = 1\), \(2\), \(\dots\), \(r\)),
    且当 \(i > r + 1\) 时, \([G']_{i,n+1} = 0\),
    且当 \(i > r\), 且 \(j \leq n\) 时,
    \([G']_{i,j} = 0\),
    且当 \(i \leq r\), 且 \(\ell \neq i\) 时,
    \([G']_{\ell,k_i} = 0\).
    (注意, 若 \(r = n\), 则 \([G']_{r+1,n+1}\) 不出现.)
    则 \(G'\) 的前 \(n\)~列作成 \(n\)~级 SREF 阵 \(A'\).
    记 \(G'\) 的列~\(n+1\) 为 \(B'\).
    则 \(G' = [A', B']\).

    为方便, 记 \([B']_{i,1} = b'_i\).
    则当 \(i > r + 1\) 时, \(b'_i = 0\).
    记 \([A']_{i,j} = a'_{i,j}\).
    则当 \(i > r\), 且 \(j \leq n\) 时,
    \(a'_{i,j} = 0\),
    且当 \(i \leq r\), 且 \(\ell \neq i\) 时,
    \(a'_{\ell,k_i} = 0\).
    则 \(G'\) 形如
    \begin{align*}
        \begin{bmatrix}
            {}         & a'_{1,k_1} & \cdots & 0          & \cdots     & 0          & \cdots & b'_1     \\
            {}         & 0          & \cdots & a'_{2,k_2} & \cdots     & 0          & \cdots & b'_2     \\
            \mathsf{O} & \vdots     & {}     & \vdots     & {}         & \vdots     & \cdots & \vdots   \\
            {}         & 0          & \cdots & 0          & \cdots     & a'_{r,k_r} & \cdots & b'_r     \\
            {}         & 0          & \cdots & 0          & \cdots     & 0          & \cdots & b'_{r+1} \\
            {}         & {}         & {}     & {}         & \mathsf{O} & {}         & {}     & {}       \\
        \end{bmatrix}.
    \end{align*}
    设从 \(1\), \(2\), \(\dots\), \(n\)
    去除 \(k_1\), \(\dots\), \(k_r\) 后,
    还剩 \(n - r\)~个数.
    我们从小到大地叫这 \(n - r\)~个数为
    \(\ell_1\), \(\dots\), \(\ell_{n-r}\).
    (注意, 若 \(r = n\), 则
    \(\ell_1\), \(\dots\), \(\ell_{n-r}\)
    不出现.)
    则 \(G'\) 对应线性方程组
    \begin{equation}
        \left\{
        \begin{aligned}
            a'_{1,k_1} x_{k_1} + \sum_{1 \leq v \leq n-r} {a'_{1,\ell_v} x_{\ell_v}} & = b'_1,     \\
            a'_{2,k_2} x_{k_2} + \sum_{1 \leq v \leq n-r} {a'_{2,\ell_v} x_{\ell_v}} & = b'_2,     \\
                                                                                     & \dots,      \\
            a'_{r,k_r} x_{k_r} + \sum_{1 \leq v \leq n-r} {a'_{r,\ell_v} x_{\ell_v}} & = b'_r,     \\
            0                                                                        & = b'_{r+1}, \\
            0                                                                        & = 0,        \\
                                                                                     & \dots,      \\
            0                                                                        & = 0,
        \end{aligned}
        \right.
        \label{eq:sole-11}
    \end{equation}
    其中, \(1 \leq k_1 < k_2 < \dots < k_r \leq n\),
    且 \(a'_{i,k_i} \neq 0\)
    (\(i = 1\), \(2\), \(\dots\), \(r\)).
    (注意, 若 \(r = n\), 则
    \begin{align*}
        + \sum_{1 \leq v \leq n-r} {a'_{i,\ell_v} x_{\ell_v}},
        \quad
        i = 1, 2, \dots, r
    \end{align*}
    不出现, 或被认为是 \(0\).)
    由前面的讨论, 方程组~\eqref{eq:sole-11}
    与方程组~\eqref{eq:sole-10} 有相同的解.

    设形如 \(E(m; p, q; s)\) (\(p \neq q\)) 的阵
    \(E_1\), \(E_2\), \(\dots\), \(E_w\) 使
    \begin{align*}
        G' = E_w (\dots (E_2 (E_1 G))).
    \end{align*}
    记 \(P = E_w \dots E_2 E_1\).
    则 \(G' = P G\).
    设 \(G\) 的列
    \(1\), \(2\), \(\dots\), \(n\), \(n+1\)
    是 \(g_1\), \(g_2\), \(\dots\), \(g_n\), \(g_{n+1}\).
    注意, \(g_j\) 是 \(A\) 的列~\(j\)
    (\(j = 1\), \(2\), \(\dots\), \(n\)),
    且 \(g_{n+1} = B\).
    则
    \begin{align*}
        G' = [P g_1, P g_2, \dots, P g_n, P g_{n+1}].
    \end{align*}
    则
    \begin{align*}
        A'
        = [P g_1, P g_2, \dots, P g_n]
        = P [g_1, g_2, \dots, g_n]
        = P A.
    \end{align*}
    则
    \begin{align*}
        A' = E_w (\dots (E_2 (E_1 A))).
    \end{align*}
    类似地, \(B' = P g_{n+1}
    = P B
    = E_w (\dots (E_2 (E_1 B)))\).

    好的.
    我们已好地准备.

    (1)
    若 \(\det {(A)} \neq 0\),
    则因为我们可用行的倍加,
    变 \(A\) 为 \(A'\),
    \(\det {(A')} = \det {(A)} \neq 0\).
    则 \(r = n\).
    则 \(k_i = i\) (\(i = 1\), \(2\), \(\dots\), \(n\)).
    则 \(G'\) 形如
    \begin{align*}
        \begin{bmatrix}
            a'_{1,1} & 0        & \cdots & 0        & b'_1   \\
            0        & a'_{2,2} & \cdots & 0        & b'_2   \\
            \vdots   & \vdots   & {}     & \vdots   & \vdots \\
            0        & 0        & \cdots & a'_{n,n} & b'_n   \\
        \end{bmatrix},
    \end{align*}
    且方程组~\eqref{eq:sole-11} 形如
    \begin{equation*}
        \left\{
        \begin{aligned}
            a'_{1,1} x_1 & = b'_1, \\
            a'_{2,2} x_2 & = b'_2, \\
                         & \dots,  \\
            a'_{n,n} x_n & = b'_n,
        \end{aligned}
        \right.
    \end{equation*}
    其中, \(a'_{i,i} \neq 0\).
    则方程组~\eqref{eq:sole-11} 有唯一的解
    \begin{equation}
        x_i = \frac{b'_i}{a'_{i,i}},
        \quad i = 1, 2, \dots, n.
        \label{eq:sole-10-solutions}
    \end{equation}
    则方程组~\eqref{eq:sole-10} 有唯一的解
    \eqref{eq:sole-10-solutions}.

    注意,
    \begin{align*}
        A' \{i, B'\}
        = {} &
        ([\dots, P g_{i-1}, P g_i, P g_{i+1}, \dots])
        \{i, P B\}
        \\
        = {} &
        [\dots, P g_{i-1}, P B, P g_{i+1}, \dots]
        \\
        = {} &
        P [\dots, g_{i-1}, B, g_{i+1}, \dots]
        \\
        = {} &
        P (A\{i, B\})
        \\
        = {} &
        E_w (\dots (E_2 (E_1 (A\{i, B\})))).
    \end{align*}
    则我们可用行的倍加,
    变 \(A\{i, B\}\) 为 \(A' \{i, B'\}\).
    则
    \begin{align*}
        \frac{\det {(A\{i, B\})}}{\det {(A)}}
        = \frac{\det {(A' \{i, B'\})}}{\det {(A')}}.
    \end{align*}
    注意, 对 \(A' \{i, B'\}\) (或 \(A'\)),
    若 \(j \neq i\),
    则 \([A' \{i, B'\}]_{i,j} = 0\)
    (或 \([A']_{i,j} = 0\)).
    则, 按行~\(i\) 展开,
    \begin{align*}
        \frac{\det {(A' \{i, B'\})}}{\det {(A')}}
        = {} &
        \frac%
        {(-1)^{i+i} [A' \{i, B'\}]_{i,i}
            \det {((A' \{i, B'\}) (i|i))}}%
        {(-1)^{i+i} [A']_{i,i}
            \det {(A' (i|i))}}
        \\
        = {} &
        \frac%
        {(-1)^{i+i} [B']_{i,1}
            \det {((A' \{i, B'\}) (i|i))}}%
        {(-1)^{i+i} [A']_{i,i}
            \det {(A' (i|i))}}
        \\
        = {} &
        \frac%
        {(-1)^{i+i} b'_i
        \det {(A' (i|i))}}%
        {(-1)^{i+i} a'_{i,i}
        \det {(A' (i|i))}}
        \\
        = {} &
        \frac{b'_i}{a'_{i,i}}
        \\
        = {} &
        x_i.
    \end{align*}
    则方程组~\eqref{eq:sole-10}
    的唯一的解~\eqref{eq:sole-10-solutions} 可被写为
    \begin{align*}
        x_i
        = \frac{\det {(A' \{i, B'\})}}{\det {(A')}}
        = \frac{\det {(A\{i, B\})}}{\det {(A)}},
        \quad
        i = 1, 2, \dots, n.
    \end{align*}

    (2)
    若 \(\det {(A)} = 0\),
    则 \(\det {(A')} = \det {(A)} = 0\).
    则 \(r < n\).

    若 \(b'_{r,n+1} \neq 0\),
    则方程组~\eqref{eq:sole-11} 无解.
    则方程组~\eqref{eq:sole-10} 无解.

    若 \(b'_{r,n+1} = 0\),
    则方程组~\eqref{eq:sole-11} 有解.
    则方程组~\eqref{eq:sole-10} 有解.
    因为 \(r < n\),
    方程组~\eqref{eq:sole-11} 有不唯一的解.
    则方程组~\eqref{eq:sole-10} 有不唯一的解.
\end{proof}

我们进一步地讨论 REF 阵.

我们知道, 我们能用行的倍加, 变阵 \(A\) 为某 REF 阵 \(B\).
不过, 变 \(A\) 为 REF 阵的方式一般是不唯一的.
设我们也用行的倍加, 变 \(A\) 为某 REF 阵 \(C\).
我们研究 \(B\) 与 \(C\) 的关系.

设用若干次行的倍加, 我们可变 \(m \times n\)~阵 \(A\)
为 REF 阵 \(B\).
则存在形如 \(E(m; p, q; s)\) (\(p \neq q\)) 的阵
\(E_1\), \(E_2\), \(\dots\), \(E_u\),
使
\begin{align*}
    B
    = E_u (E_{u-1} \dots (E_2 (E_1 A)))
    = (E_u E_{u-1} \dots E_2 E_1) A.
\end{align*}
类似地,
设用若干次行的倍加, 我们可变 \(m \times n\)~阵 \(A\)
为 REF 阵 \(C\).
则存在形如 \(E(m; p, q; s)\) (\(p \neq q\)) 的阵
\(F_1\), \(F_2\), \(\dots\), \(F_v\),
使
\begin{align*}
    C
    = F_v (F_{v-1} \dots (F_2 (F_1 A)))
    = (F_v F_{v-1} \dots F_2 F_1) A.
\end{align*}
我们知道, 对每个 \(E_i\) (或 \(F_j\)),
存在形如 \(E(m; p, q; s)\) (\(p \neq q\)) 的阵
\(G_i\) (或 \(H_j\)),
使 \(G_i E_i = I_m\)
(或 \(H_j F_j = I_m\)).
则
\begin{align*}
    (G_1 G_2 \dots G_{u-1} G_u) B
    = {} &
    (G_1 G_2 \dots G_{u-1} G_u)
    (E_u E_{u-1} \dots E_2 E_1) A
    \\
    = {} &
    (G_1 (G_2 (\dots (G_{u-1} (G_u E_u) E_{u-1}) \dots) E_2) E_1)
    A
    \\
    = {} &
    I_m A
    \\
    = {} &
    A.
\end{align*}
类似地, \((H_1 H_2 \dots H_{v-1} H_v) C = A\).
则
\begin{align*}
    B
    = (E_u \dots E_1) A
    = (E_u \dots E_1) ((H_1 \dots H_v) C)
    = ((E_u \dots E_1) (H_1 \dots H_v)) C.
\end{align*}
则存在 \(m\)~级阵 \(P = (E_u \dots E_1) (H_1 \dots H_v)\),
使 \(B = P C\).
类似地, 也存在 \(m\)~级阵
\(Q = (F_v \dots F_1) (G_1 \dots G_u)\),
使 \(C = Q B\).

由此, 我们可以证明:

\begin{theorem}
    设 \(B\), \(C\) 是 \(m \times n\) REF 阵.
    记 \(k_i = \operatorname{pivot} {(B; i)}\);
    记 \(k'_i = \operatorname{pivot} {(C; i)}\).
    设 \(k_1 < k_2 < \dots < k_r < \infty\),
    且对 \(i > r\), \(B\) 的行~\(i\) 是 \(0\).
    设 \(k'_1 \allowbreak
    < k'_2 \allowbreak
    < \dots \allowbreak
    < k'_{r'} \allowbreak
    < \infty\),
    且对 \(i' > r'\), \(C\) 的行~\(i'\) 是 \(0\).
    设 \(m\)~级阵 \(P\), \(Q\) 适合,
    \(B = PC\), 且 \(C = QB\).
    则 \(r = r'\),
    且对 \(i \leq r\),
    \(k_i = k'_i\).
\end{theorem}

\begin{proof}
    我们证明, \(k_1 = k'_1\).
    为此, 我们先证明, \(k_1 \geq k'_1\).
    反设 \(k_1 < k'_1\).
    \(C\) 是 REF 阵, 且 \(k_1 < k'_1\),
    故若 \(v \geq 1\), 则 \([C]_{v,k_1} = 0\).
    \(B\) 是 REF 阵, 故 \([B]_{1,k_1} \neq 0\).
    但
    \begin{align*}
        [B]_{1,k_1}
        = {} &
        [P C]_{1,k_1}
        \\
        = {} &
        \sum_{1 \leq v \leq m} {[P]_{1,v} [C]_{v,k_1}}
        \\
        = {} &
        \sum_{1 \leq v \leq m} {[P]_{1,v} 0}
        \\
        = {} &
        0.
    \end{align*}
    这是矛盾.
    故 \(k_1 \geq k'_1\).
    类似地, 由 \(C = Q P\),
    我们可证明, \(k_1 \leq k'_1\).
    则 \(k_1 = k'_1\).

    设我们已知 \(k_\ell = k'_\ell\)
    (\(\ell \leq t\), 且 \(t < r\), 且 \(t < r'\)).
    我们由此证明, \(k_{t+1} = k'_{t+1}\).
    为此, 我们先证明, \(k_{t+1} \geq k'_{t+1}\).
    反设 \(k_{t+1} < k'_{t+1}\).
    \(B\) 是 REF 阵,
    故对 \(\ell \leq t\),
    \([B]_{t+1,k_\ell} = 0\).
    则
    \begin{align*}
        0
        = {} &
        [B]_{t+1,k_\ell}
        \\
        = {} &
        [P C]_{t+1,k_\ell}
        \\
        = {} &
        \sum_{1 \leq v \leq m}
        {[P]_{t+1,v} [C]_{v,k_\ell}}
        \\
        = {} &
        \sum_{1 \leq v \leq m}
        {[P]_{t+1,v} [C]_{v,k'_\ell}}
        \\
        = {} &
        \sum_{1 \leq v \leq \ell}
        {[P]_{t+1,v} [C]_{v,k'_\ell}}
        +
        \sum_{\ell < v \leq m}
        {[P]_{t+1,v} [C]_{v,k'_\ell}}
        \\
        = {} &
        \sum_{1 \leq v \leq \ell}
        {[P]_{t+1,v} [C]_{v,k'_\ell}}
        +
        \sum_{\ell < v \leq m}
        {[P]_{t+1,v} 0}
        \\
        = {} &
        \sum_{1 \leq v \leq \ell}
        {[C]_{v,k'_\ell} [P]_{t+1,v}}.
    \end{align*}
    则
    \begin{equation*}
        \begin{alignedat}{7}
            [C]_{1,k'_1} [P]_{t+1,1} &         &                          &         &       &         &                          & {} = 0, \\
            [C]_{1,k'_2} [P]_{t+1,1} & {} + {} & [C]_{2,k'_2} [P]_{t+1,2} &         &       &         &                          & {} = 0, \\
                                     &         &                          &         &       &         &                          & \dots,  \\
            [C]_{1,k'_t} [P]_{t+1,1} & {} + {} & [C]_{2,k'_t} [P]_{t+1,2} & {} + {} & \dots & {} + {} & [C]_{t,k'_t} [P]_{t+1,t} & {} = 0.
        \end{alignedat}
    \end{equation*}
    因为 \([C]_{\ell,k'_\ell} \neq 0\),
    我们可从前向后地确定,
    \([P]_{t+1,1}\), \([P]_{t+1,2}\), \(\dots\), \([P]_{t+1,t}\)
    都是 \(0\).
    注意, \(C\) 是 REF 阵,
    且 \(k_{t+1} < k'_{t+1}\),
    故若 \(v \geq t+1\),
    则 \([C]_{v,k_{t+1}} = 0\).
    注意, \([B]_{t+1,k_{t+1}} \neq 0\).
    但
    \begin{align*}
        [B]_{t+1,k_{t+1}}
        = {} &
        [P C]_{t+1,k_{t+1}}
        \\
        = {} &
        \sum_{1 \leq v \leq m}
        {[P]_{t+1,v} [C]_{v,k_{t+1}}}
        \\
        = {} &
        \sum_{1 \leq v \leq t}
        {[P]_{t+1,v} [C]_{v,k_{t+1}}}
        +
        \sum_{t < v \leq m}
        {[P]_{t+1,v} [C]_{v,k_{t+1}}}
        \\
        = {} &
        \sum_{1 \leq v \leq t}
        {0\, [C]_{v,k_{t+1}}}
        +
        \sum_{t < v \leq m}
        {[P]_{t+1,v} 0}
        \\
        = {} &
        0.
    \end{align*}
    这是矛盾.
    故 \(k_{t+1} \geq k'_{t+1}\).
    类似地, 由 \(C = Q P\),
    我们可证明, \(k_{t+1} \leq k'_{t+1}\).
    则 \(k_{t+1} = k'_{t+1}\).

    最后, 我们证明, \(r = r'\).
    为此, 我们先证明, \(r' \geq r\).
    反设 \(r' < r\).
    那么, 由前面的讨论, 对 \(\ell \leq r' < r\),
    \(k_\ell = k'_\ell\).
    \(B\) 是 REF 阵,
    故对 \(\ell \leq r' < r\),
    \([B]_{r,k_\ell} = 0\).
    注意, 对 \(v > r'\),
    \([C]_{v,j} = 0\).
    则
    \begin{align*}
        0
        = {} &
        [B]_{r,k_\ell}
        \\
        = {} &
        [P C]_{r,k_\ell}
        \\
        = {} &
        \sum_{1 \leq v \leq m}
        {[P]_{r,v} [C]_{v,k_\ell}}
        \\
        = {} &
        \sum_{1 \leq v \leq r'}
        {[P]_{r,v} [C]_{v,k_\ell}}
        +
        \sum_{r' < v \leq m}
        {[P]_{r,v} [C]_{v,k_\ell}}
        \\
        = {} &
        \sum_{1 \leq v \leq r'}
        {[P]_{r,v} [C]_{v,k'_\ell}}
        +
        \sum_{r' < v \leq m}
        {[P]_{r,v} 0}
        \\
        = {} &
        \sum_{1 \leq v \leq r'}
        {[C]_{v,k'_\ell} [P]_{r,v}}
        \\
        = {} &
        \sum_{1 \leq v \leq \ell}
        {[C]_{v,k'_\ell} [P]_{r,v}}
        +
        \sum_{\ell < v \leq r'}
        {[C]_{v,k'_\ell} [P]_{r,v}}
        \\
        = {} &
        \sum_{1 \leq v \leq \ell}
        {[C]_{v,k'_\ell} [P]_{r,v}}
        +
        \sum_{\ell < v \leq r'}
        {0\, [P]_{r,v}}
        \\
        = {} &
        \sum_{1 \leq v \leq \ell}
        {[C]_{v,k'_\ell} [P]_{r,v}}.
    \end{align*}
    则
    \begin{equation*}
        \begin{alignedat}{7}
            [C]_{1,k'_1} [P]_{r,1}    &         &                           &         &       &         &                             & {} = 0, \\
            [C]_{1,k'_2} [P]_{r,1}    & {} + {} & [C]_{2,k'_2} [P]_{r,2}    &         &       &         &                             & {} = 0, \\
                                      &         &                           &         &       &         &                             & \dots,  \\
            [C]_{1,k'_{r'}} [P]_{r,1} & {} + {} & [C]_{2,k'_{r'}} [P]_{r,2} & {} + {} & \dots & {} + {} & [C]_{r',k'_{r'}} [P]_{r,r'} & {} = 0.
        \end{alignedat}
    \end{equation*}
    因为 \([C]_{\ell,k'_\ell} \neq 0\),
    我们可从前向后地确定,
    \([P]_{r,1}\), \([P]_{r,2}\), \(\dots\), \([P]_{r,r'}\)
    都是 \(0\).
    则
    \begin{align*}
        [B]_{r,j}
        = {} &
        [P C]_{r,j}
        \\
        = {} &
        \sum_{1 \leq v \leq m}
        {[P]_{r,v} [C]_{v,j}}
        \\
        = {} &
        \sum_{1 \leq v \leq r'}
        {[P]_{r,v} [C]_{v,j}}
        +
        \sum_{r' < v \leq m}
        {[P]_{r,v} [C]_{v,j}}
        \\
        = {} &
        \sum_{1 \leq v \leq r'}
        {0\, [C]_{v,j}}
        +
        \sum_{r' < v \leq m}
        {[P]_{r,v} 0}
        \\
        = {} &
        0.
    \end{align*}
    则 \(B\) 的行~\(r\) 是 \(0\).
    这是矛盾, 因为 \([B]_{r,k_r} \neq 0\).
    故 \(r' \geq r\).
    类似地, 由 \(C = Q P\),
    我们可证明, \(r \geq r'\).
    则 \(r = r'\).
\end{proof}

值得提, 我们也能如此证明, \(r = r'\).
注意,
\(\operatorname{rank} {(B)}
= \operatorname{rank} {(P C)}
\leq \operatorname{rank} {(C)}\),
且
\(\operatorname{rank} {(C)}
= \operatorname{rank} {(Q B)}
\leq \operatorname{rank} {(B)}\),
故 \(\operatorname{rank} {(B)}
= \operatorname{rank} {(C)}\).
回想, 对 REF 阵 \(B\) 与 \(C\),
\(B\) 有 \(r = \operatorname{rank} {(B)}\) 个非零行,
且 \(C\) 有 \(r' = \operatorname{rank} {(C)}\) 个非零行.
则 \(r = r'\).
由此可见, 秩 (与行列式) 是有用的工具.

进一步地, 对 SREF 阵, 我们有:

\begin{theorem}
    设 \(B\), \(C\) 是 \(m \times n\) SREF 阵.
    记 \(k_i = \operatorname{pivot} {(B; i)}\);
    记 \(k'_i = \operatorname{pivot} {(C; i)}\).
    设 \(k_1 < k_2 < \dots < k_r < \infty\),
    且对 \(i > r\), \(B\) 的行~\(i\) 是 \(0\).
    设 \(k'_1 \allowbreak
    < k'_2 \allowbreak
    < \dots \allowbreak
    < k'_{r'} \allowbreak
    < \infty\),
    且对 \(i' > r'\), \(C\) 的行~\(i'\) 是 \(0\).
    设 \(m\)~级阵 \(P\), \(Q\) 适合,
    \(B = PC\), 且 \(C = QB\).
    则 \(r = r'\),
    且对 \(i \leq r\),
    \(k_i = k'_i\),
    且对 \(i \leq r\),
    \(B\) 的行~\(i\) 是 \(C\) 的行~\(i\) 的
    \([B]_{i,k_i} / [C]_{i,k_i}\)~倍.
\end{theorem}

\begin{proof}
    SREF 阵是 REF 阵,
    故由上个定理, 我们再确定元的关系即可.

    注意, \(C\) 是 SREF 阵.
    则对 \(i\), \(\ell \leq r\),
    \begin{align*}
        [B]_{i,k_\ell}
        = {} &
        [P C]_{i,k_\ell}
        \\
        = {} &
        \sum_{1 \leq v \leq m}
        {[P]_{i,v} [C]_{v,k_\ell}}
        \\
        = {} &
        [P]_{i,\ell} [C]_{\ell,k_\ell}
        +
        \sum_{\substack{1 \leq v \leq m \\ v \neq \ell}}
        {[P]_{i,v} [C]_{v,k_\ell}}
        \\
        = {} &
        [P]_{i,\ell} [C]_{\ell,k_\ell}
        +
        \sum_{\substack{1 \leq v \leq m \\ v \neq \ell}}
        {[P]_{i,v} 0}
        \\
        = {} &
        [P]_{i,\ell} [C]_{\ell,k_\ell}.
    \end{align*}
    当 \(i \neq \ell\) 时, \([B]_{i,k_\ell} = 0\).
    则当 \(i \neq \ell\) 时, \([P]_{i,\ell} = 0\).
    另一方面, 我们取 \(\ell = i\),
    得 \([P]_{i,i} = [B]_{i,k_i} / [C]_{i,k_i}\).
    则当 \(i \leq r\) 时,
    \begin{align*}
        [B]_{i,j}
        = {} &
        [P C]_{i,j}
        \\
        = {} &
        \sum_{1 \leq v \leq m}
        {[P]_{i,v} [C]_{v,j}}
        \\
        = {} &
        [P]_{i,i} [C]_{i,j}
        +
        \sum_{\substack{1 \leq v \leq m \\ v \neq i}}
        {[P]_{i,v} [C]_{v,j}}
        \\
        = {} &
        \frac{[B]_{i,k_i}}{[C]_{i,k_i}}
        [C]_{i,j}
        +
        \sum_{\substack{1 \leq v \leq m \\ v \neq i}}
        {0\, [C]_{v,j}}
        \\
        = {} &
        \frac{[B]_{i,k_i}}{[C]_{i,k_i}}
        [C]_{i,j}.
        \qedhere
    \end{align*}
\end{proof}

回想, 我们定义了 (列的) 倍乘:
以某数乘阵的一列.
类似地, 我们可考虑行的倍乘:
以数 \(s\) 乘 \(m \times n\)~阵 \(A\) 的行~\(q\),
且不改变其他的行, 得阵 \(B\).
则我们可类似地证明,
\(B = M(m; q; s) A\).

若 \(s \neq 0\), 我们可以证明,
\(M(m; q; s^{-1}) M(m; q; s) = I_m\):
考虑行~\(q\) 的元即可.
类似地, 若 \(s \neq 0\),
\(M(m; q; s) M(m; q; s^{-1}) = I_m\).

若 \(s \neq 0\), 则对 \(m \times n\)~阵 \(A\),
\begin{align*}
    \operatorname{rank} {(A)}
    \geq {} &
    \operatorname{rank} {(M(m; q; s) A)}
    \\
    \geq {} &
    \operatorname{rank}
    {(M(m; q; s^{-1}) (M(m; q; s) A))}
    \\
    = {}    &
    \operatorname{rank}
    {((M(m; q; s^{-1}) M(m; q; s)) A)}
    \\
    = {}    &
    \operatorname{rank} {(I_m A)}
    \\
    = {}    &
    \operatorname{rank} {(A)}.
\end{align*}
故 \(\operatorname{rank} {(A)}
\geq \operatorname{rank} {(M(m; q; s) A)}
\geq \operatorname{rank} {(A)}\).
于是, 以非零的数乘阵的一行不改变秩.
(类似地, 以非零的数乘阵的一列不改变秩.)

% 类似地, 我们可以证明,
% 以非零的数乘线性方程组的一个方程不改变它的解:
% 若 \(a_{i,1} c_1 + a_{i,2} c_2 + \dots + a_{i,n} c_n = b_i\),
% 则
% \begin{align*}
%          &
%     (s a_{i,1}) c_1 + (s a_{i,2}) c_2 + \dots + (s a_{i,n}) c_n
%     \\
%     = {} &
%     s (a_{i,1} c_1) + s (a_{i,2} c_2) + \dots + s (a_{i,n} c_n)
%     \\
%     = {} &
%     s (a_{i,1} c_1 + a_{i,2} c_2 + \dots + a_{i,n} c_n)
%     \\
%     = {} &
%     s b_i;
% \end{align*}
% 反过来, 若
% \((s a_{i,1}) d_1 + (s a_{i,2}) d_2 + \dots + (s a_{i,n}) d_n
% = s b_i\),
% 则因 \(s^{-1} s = 1\),
% \begin{align*}
%          &
%     a_{i,1} d_1 + a_{i,2} d_2 + \dots + a_{i,n} d_n
%     \\
%     = {} &
%     (s^{-1} s) a_{i,1} d_1 + (s^{-1} s) a_{i,2} d_2
%     + \dots + (s^{-1} s) a_{i,n} d_n
%     \\
%     = {} &
%     s^{-1} (s a_{i,1}) d_1 + s^{-1} (s a_{i,2}) d_2
%     + \dots + s^{-1} (s a_{i,n}) d_n
%     \\
%     = {} &
%     s^{-1}
%     ((s a_{i,1}) d_1 + (s a_{i,2}) d_2 + \dots + (s a_{i,n}) d_n)
%     \\
%     = {} &
%     s^{-1} (s b_i)
%     \\
%     = {} &
%     (s^{-1} s) b_i
%     \\
%     = {} &
%     b_i.
% \end{align*}

类似地, 我们可以证明,
以非零的数乘线性方程组的一个方程不改变它的解:
当 \(s \neq 0\) 时,
若
\(a_{q,1} c_1 + a_{q,2} c_2 + \dots + a_{q,n} c_n = b_q\),
则
\((s a_{q,1}) c_1 + (s a_{q,2}) c_2 + \dots + (s a_{q,n}) c_n
= s b_q\),
且若
\((s a_{q,1}) d_1 + (s a_{q,2}) d_2 + \dots + (s a_{q,n}) d_n
= s b_q\),
则
\(a_{q,1} d_1 + a_{q,2} d_2 + \dots + a_{q,n} d_n = b_q\).

\begin{definition}
    设 \(A\) 是 \(m \times n\)~阵.
    设存在不高于 \(m\) 的非负整数 \(r\), 使:

    (1)
    对任何不高于 \(r\) 的正整数 \(i\),
    \begin{align*}
        \operatorname{pivot} {(A; i)} < \infty;
    \end{align*}

    (2)
    对任何低于 \(r\) 的正整数 \(i\),
    \begin{align*}
        \operatorname{pivot} {(A; i)}
        < \operatorname{pivot} {(A; i + 1)};
    \end{align*}

    (3)
    对任何高于 \(r\) 且不高于 \(m\) 的整数 \(i\),
    \begin{align*}
        \operatorname{pivot} {(A; i)} = \infty;
    \end{align*}

    (4)
    对任何不高于 \(r\) 的正整数 \(i\),
    与任何不高于 \(m\) 的正整数 \(\ell\),
    \begin{align*}
        [A]_{\ell,\operatorname{pivot} {(A; i)}} =
        \begin{cases}
            0, & \ell \neq i; \\
            1, & \ell = i.
        \end{cases}
    \end{align*}

    我们说, \(A\) 是 \emph{RREF}
    (\angla{reduced row-echelon form}) \emph{阵}.
\end{definition}

注意, RREF 阵是 SREF 阵,
但 SREF 阵不一定是 RREF 阵.
比如, \([-1]\) 是 SREF 阵, 但它不是 RREF 阵.

我们知道, 我们可用行的倍加,
变一个阵为 SREF 阵.
类似地,

\begin{theorem}
    设 \(A\) 是 \(m \times n\)~阵.
    则我们可用行的倍加与行的倍乘
    (且不以 \(0\) 乘阵的行),
    变 \(A\) 为某 \(m \times n\) RREF 阵 \(D\).
\end{theorem}

\begin{proof}
    我们可用行的倍加,
    变 \(A\) 为某 \(m \times n\) SREF 阵 \(B\).
    则存在非负整数 \(r\), 使
    \begin{align*}
        \operatorname{pivot} {(B; 1)}
        < \operatorname{pivot} {(B; 2)}
        < \dots
        < \operatorname{pivot} {(B; r-1)}
        < \operatorname{pivot} {(B; r)}
        < \infty,
    \end{align*}
    且对 \(i > r\),
    \(B\) 的行~\(i\) 是 \(0\),
    且若 \(\ell \neq i\),
    则 \([B]_{\ell,\operatorname{pivot} {(B; i)}} = 0\).
    对 \(B\),
    我们分别地以
    \(1/[B]_{i,\operatorname{pivot} {(B; i)}} \neq 0\)
    乘行~\(i\)
    (\(i = 1, 2, \dots, r\)),
    得阵 \(D\).
    则 \(\operatorname{pivot} {(D; i)}
    = \operatorname{pivot} {(B; i)}\),
    且 \([D]_{i,\operatorname{pivot} {(D; i)}} = 1\).
    则 \(D\) 是 RREF 阵.
\end{proof}

\begin{example}
    设 SREF 阵
    \begin{align*}
        G_2 =
        \begin{bmatrix}
            2 & 0  & 46  \\
            0 & -1 & -12 \\
        \end{bmatrix}.
    \end{align*}
    它不是 RREF 阵.
    不过, 对 \(G_2\),
    我们分别地以 \(1/2\), \(-1\) 乘行~\(1\), \(2\),
    得 RREF 阵
    \begin{align*}
        G_3 =
        \begin{bmatrix}
            1 & 0 & 23 \\
            0 & 1 & 12 \\
        \end{bmatrix}.
    \end{align*}
\end{example}

注意, 只作行的倍加不一定能变阵为 RREF 阵.

\begin{example}
    设 SREF 阵
    \begin{align*}
        A =
        \begin{bmatrix}
            1 & 0  \\
            0 & -1 \\
        \end{bmatrix}.
    \end{align*}
    我们说, 我们不能只作行的倍加,
    以变 \(A\) 为 RREF 阵.

    反设, 我们能只作行的倍加,
    以变 \(A\) 为 RREF 阵 \(B\).
    则存在 \(2\)~级阵 \(P\), \(Q\),
    使 \(B = P A\), 且 \(A = Q B\).
    注意, \(A\) 与 \(B\) 都是 SREF 阵.
    则由前面的讨论, \(B\) 应形如
    \begin{align*}
        \begin{bmatrix}
            [B]_{1,1} & 0         \\
            0         & [B]_{2,2} \\
        \end{bmatrix},
    \end{align*}
    其中, \([B]_{i,i} \neq 0\).
    \(B\) 是 RREF 阵, 故
    \([B]_{1,1} = [B]_{2,2} = 1\).
    则 \(\det {(B)} = 1\).
    注意, 倍加不改变行列式,
    但 \(\det {(A)} = -1 \neq \det {(B)}\).
    这是矛盾.
\end{example}

由前面的讨论, 对 RREF 阵, 我们有:

\begin{theorem}
    设 \(B\), \(B'\) 是 \(m \times n\) RREF 阵.
    设 \(m\)~级阵 \(P\), \(Q\) 适合,
    \(B = P B'\), 且 \(B' = Q B\).
    则 \(B = B'\).
\end{theorem}

\begin{proof}
    RREF 阵是 SREF 阵.
    则 \(\operatorname{pivot} {(B; i)}
    = \operatorname{pivot} {(B'; i)}\),
    且 \(B\), \(B'\) 的非零行的数目相等
    (其我们设是 \(r\)),
    且对 \(i \leq r\),
    \(B\) 的行~\(i\) 是 \(B'\) 的行~\(i\) 的
    \(1/1 = 1\)~倍
    (注意, 对 \(i \leq r\),
    \([B]_{i,\operatorname{pivot} {(B; i)}}
    = 1 = [B']_{i,\operatorname{pivot} {(B'; i)}}\)).
    对 \(i > r\),
    \(B\) 的行~\(i\) 与 \(B'\) 的行~\(i\) 都是 \(0\).
    则 \(B = B'\).
\end{proof}

设 \(A\) 是 \(m \times n\)~阵.
设我们用行的倍加与行的倍乘
(且不以 \(0\) 乘阵的行, 下同),
变 \(A\) 为某 \(m \times n\) RREF 阵 \(B\).
设我们用行的倍加与行的倍乘,
变 \(A\) 为某 \(m \times n\) RREF 阵 \(B'\).
则存在形如 \(E(m; p, q; s)\) (\(p \neq q\))
或 \(M(m; q; s)\) (\(s \neq 0\)) 的阵
\(C_1\), \(C_2\), \(\dots\), \(C_u\),
与 \(D_1\), \(D_2\), \(\dots\), \(D_v\),
使
\begin{align*}
    B  & = C_u (C_{u-1} \dots (C_2 (C_1 A)))
    = (C_u C_{u-1} \dots C_2 C_1) A,
    \\
    B' & = D_v (D_{v-1} \dots (D_2 (D_1 A)))
    = (D_v D_{v-1} \dots D_2 D_1) A.
\end{align*}
对每个 \(C_i\) (或 \(D_j\)),
存在形如 \(E(m; p, q; s)\) (\(p \neq q\))
或 \(M(m; q; s)\) (\(s \neq 0\)) 的阵
\(F_i\) (或 \(G_j\)),
使 \(F_i C_i = I_m\)
(或 \(G_j D_j = I_m\)).
则, 可以验证,
\begin{align*}
    (F_1 F_2 \dots F_{u-1} F_u) B
    = A
    = (G_1 G_2 \dots G_{v-1} G_v) B'.
\end{align*}
则
\begin{align*}
    B  & = (C_u \dots C_1) A
    = ((C_u \dots C_1) (G_1 \dots G_v)) B',
    \\
    B' & = (D_v \dots D_1) A
    = ((D_v \dots D_1) (F_1 \dots F_u)) B.
\end{align*}
则存在 \(m\)~级阵 \(P\), \(Q\),
使 \(B = P B'\), 且 \(B' = Q B\).
我们有:

\begin{theorem}
    设 \(A\) 是 \(m \times n\)~阵.
    设我们用行的倍加与行的倍乘
    (且不以 \(0\) 乘阵的行, 下同),
    变 \(A\) 为某 \(m \times n\) RREF 阵 \(B\).
    设我们用行的倍加与行的倍乘,
    变 \(A\) 为某 \(m \times n\) RREF 阵 \(B'\).
    则 \(B = B'\).
\end{theorem}

\end{document}

This material discusses the rank of a matrix,
matrices in row-echelon form,
and systems of linear equations.
